% ============================================================
% CHAPTER 05
% IDENTITY AND CONTINUITY
% Book I: Ontology of Viability
% Part I: Identity and Manifestation
% ============================================================

\chapter{Identity and Continuity}
\label{ch:identity-and-continuity}

% ------------------------------------------------------------
% CHAPTER METADATA
% ------------------------------------------------------------

\newcommand{\ChapterFiveTitle}{Identity and Continuity}
\newcommand{\ChapterFiveID}{VS.B01.P01.C05}

% ------------------------------------------------------------
% CHAPTER OPENING
% ------------------------------------------------------------

The relation between identity and continuity constitutes
a fundamental problem in the ontology of an existent.

An existent may undergo temporal succession, qualitative
change, structural transformation, alteration of its
manifestations, or interruption of particular processes.
Yet none of these conditions, considered independently,
is sufficient to establish whether its identity persists,
continues, or ceases to be.

The problem therefore requires a distinction between
the identity of an existent and the various forms of
continuity that may characterize its existence,
determinations, manifestations, and relations.

Identity must not be reduced to persistence, and persistence
must not be treated as equivalent to continuity.
Likewise, the continuity of an existent does not necessarily
establish the continuity of its identity, while the
discontinuity of a particular manifestation or process
does not, by itself, establish the termination of identity.

The present chapter examines these distinctions and
investigates the ontological conditions under which
identity may be preserved, continued, interrupted,
or brought into question.

It does not assume that identity is necessarily preserved
through every temporal transformation. Nor does it assume
that every significant transformation establishes the
emergence of a numerically different identity.

The inquiry remains independent of recognition, memory,
observation, and epistemic identification. Its concern
is the ontological status of identity and the conditions
under which its continuity can be established.

The chapter proceeds through seven sections, moving from
the formulation of the ontological problem to the distinction
between persistence and continuity, the relation between
invariance and temporal succession, and the examination
of the conditions and limits of identity preservation.

% ------------------------------------------------------------
% CHAPTER SECTIONS
% ------------------------------------------------------------

% =============================================================================
% Section 01 — The Problem of Identity Continuity
% Book I — Ontology of Viability
% Chapter 05 — Identity and Continuity
%
% Purpose:
%   To formulate the ontological problem of identity continuity
%   and distinguish it from persistence, temporal succession,
%   invariance, and the continuity of manifestation.
%
% Conceptual Policy:
%   1. Identity, persistence, and continuity are distinct concepts.
%   2. The continuity of an existent does not automatically
%      establish the continuity of its identity.
%   3. Difference between successive conditions does not
%      automatically establish numerical non-identity.
%   4. Invariance is neither assumed to be sufficient nor
%      necessary for identity continuity.
%   5. Neither unity nor plurality of identity is presupposed.
%   6. The analysis remains independent of observers,
%      memory, recognition, and epistemic accessibility.
% =============================================================================

\section{The Problem of Identity Continuity}
\label{sec:problem-of-identity-continuity}

\subsection{The Ontological Problem of Identity Continuity}

The question of identity continuity concerns one of the
most fundamental problems in the ontology of an existent:
under what conditions can an identity be considered
continuous through change, succession, or transformation?

An existent may undergo alterations in its determinations,
relations, composition, material constitution, or
manifestations. Such alterations raise the question
of whether the identity associated with that existent
is preserved, transformed, interrupted, or replaced.

The existence of change, however, does not by itself
settle the question of identity continuity.

Nor does the persistence of an existent automatically
establish that its identity remains the same throughout
the relevant interval.

The central problem is therefore to distinguish the
continuity of an existent from the continuity of its
identity, while examining the conditions under which
these two forms of continuity may correspond.

This inquiry must not begin by assuming that every
existent possesses exactly one identity, that identity
necessarily persists through every transformation,
or that any particular feature provides an exhaustive
criterion of identity preservation.

Identity continuity must be investigated as an
independent ontological problem.

\subsection{The Difference Between Identity and Continuity}

Identity and continuity are closely related concepts,
but their relation does not establish their equivalence.

Identity concerns the ontological question of what
makes an existent the identity that it is. Continuity
concerns the possible relation of an existent or
identity across different conditions, stages, or
configurations.

The first question concerns identity as such. The
second concerns whether and how an identity may
remain continuous across a relevant difference
or transition.

Consequently, continuity cannot be treated as the
definition of identity.

An account of continuity presupposes that the terms
being related have already been specified. It must
therefore establish whether the relevant terms are
successive conditions of one identity, conditions
associated with distinct identities, or configurations
whose numerical relation remains undetermined.

Likewise, the identification of an identity does
not by itself explain the conditions under which
that identity is continuous.

The distinction is necessary to prevent the concept
of identity from being reduced to a temporal or
relational property.

\subsection{Identity Continuity and the Continuity of an Existent}

The continuity of an existent and the continuity
of its identity must be distinguished.

An existent may persist through a sequence of
conditions, yet the question of whether the same
identity is maintained throughout that sequence
may remain unresolved.

Conversely, the continuity of an identity may not
be adequately described by considering only the
persistence of a particular material configuration
or the uninterrupted presence of a particular
manifestation.

The distinction requires an independent examination
of the relation between existence and identity.

The continued existence of something does not
automatically establish the numerical continuity
of its identity. Similarly, a change in the mode,
condition, or configuration of existence does not
automatically establish that identity has ceased
or been replaced.

The relation between existential persistence and
identity continuity must therefore be established
rather than presumed.

\subsection{The Distinction Between Continuity and Persistence}

Persistence and continuity are frequently used
to describe the continued presence of an existent.
Nevertheless, they address different analytical
questions.

Persistence concerns whether an existent or
ontological condition remains present across a
specified interval or through a sequence of changes.

Continuity concerns whether there is an appropriate
relation connecting the relevant conditions or
stages without an ontologically decisive interruption.

These preliminary distinctions do not yet establish
a complete definition of either concept. They
indicate that the continued presence of an existent
and the continuity of its successive conditions
cannot be assumed to be identical.

An existent may persist while undergoing discontinuous
changes in a particular respect. Conversely, a
continuous sequence of conditions may be associated
with a process whose persistence, identity, or
numerical unity requires independent examination.

Persistence therefore cannot serve as an unrestricted
criterion of continuity, and continuity cannot
automatically be inferred from the mere continued
existence of something.

\subsection{Identity Continuity and Temporal Succession}

Temporal succession introduces a further dimension
to the problem.

An existent may be considered at different temporal
positions or under successive temporal conditions.
The question then arises whether the identity
associated with those conditions remains numerically
the same.

However, temporal succession and identity continuity
are not identical.

The fact that one condition precedes another does
not establish that both conditions belong to the
same identity. Temporal proximity, uninterrupted
succession, or a continuous sequence of changes
cannot independently settle the numerical relation
between the relevant identities.

Likewise, the existence of a temporal interval
does not by itself establish that identity has
been interrupted.

The analysis must therefore distinguish temporal
ordering from ontological continuity.

Time may provide a domain within which identity
continuity is investigated, but temporal succession
alone does not provide a sufficient criterion
for determining it.

\subsection{The Problem of Change and Identity Preservation}

Change is central to the problem of identity continuity
because the conditions associated with an existent
may differ across successive configurations.

Such differences may concern determinations,
relations, composition, material constitution,
structural organization, or manifestation.

The existence of these differences raises the
question of whether identity can remain continuous
without requiring the preservation of every feature
or condition.

If identity continuity required the complete absence
of change, it would be reduced to immutability.
Such a reduction would exclude the possibility
of identity preservation through transformation
without establishing that this exclusion is
ontologically justified.

At the same time, the existence of change cannot
be treated as irrelevant to identity.

Some transformations may be compatible with
identity preservation, whereas others may alter
or interrupt conditions relevant to identity.
The distinction between these possibilities
requires further investigation.

The problem is therefore not whether change
occurs, but which changes, under which conditions,
are compatible with the continuity of identity.

\subsection{The Difference Between Qualitative Change and Numerical Non-Identity}

A fundamental difficulty arises when qualitative
difference is treated as evidence of numerical
non-identity.

Two conditions of an existent may differ in
their determinations, configurations, relations,
or material characteristics without thereby
establishing that they belong to numerically
distinct identities.

Qualitative difference concerns how two conditions
differ in their characteristics or organization.
Numerical non-identity concerns whether the
relevant terms are distinct identities rather
than different conditions of one identity.

These are separate questions.

The difference between two successive conditions
may be compatible with identity continuity.
However, this compatibility cannot be generalized
into the claim that every difference is compatible
with the preservation of identity.

The ontological significance of a difference
depends upon the subject, respect, and conditions
under consideration.

Accordingly, neither qualitative similarity
nor qualitative difference is sufficient by
itself to establish numerical identity or
numerical non-identity.

\subsection{Identity Continuity and Invariance}

The relation between identity continuity and
invariance requires particular care.

An existent may preserve certain features,
relations, or structural properties through
a specified transformation. Such preservation
may provide relevant information about the
continuity of the existent or one of its
ontological aspects.

Nevertheless, the presence of an invariant
does not independently establish that identity
has been preserved.

An invariant may belong to a particular
determination, structure, relation, or domain
without exhausting the identity of the existent.

Conversely, the absence of an identified
invariant does not automatically establish
that identity continuity is impossible.

Invariance must therefore be examined in
relation to specified transformations,
domains, and respects.

Identity continuity cannot be reduced to
the preservation of invariants, and the
question of whether particular invariants
are necessary for identity preservation
must remain open.

\subsection{Identity Continuity and Manifestation}

The continuity of identity must also be
distinguished from the continuity of manifestation.

An existent may undergo changes in the
manifestation of its determinations or
other ontological aspects. These changes
may occur without establishing whether
identity has been preserved or interrupted.

Likewise, the continuity of a manifestation
does not necessarily establish the continuity
of the identity associated with it.

A particular manifestation may persist
while the numerical relation of the relevant
identity remains unresolved. Multiple
manifestations may also be associated with
one identity or with distinct identities.

The relation between identity and manifestation
must therefore be investigated independently
of the temporal continuity of any particular
manifest condition.

This distinction also prevents the continuity
of an observable or materially accessible
configuration from becoming an unrestricted
criterion of identity.

\subsection{Identity Continuity and the Question of Unity}

The problem of identity continuity cannot
be separated from the question of unity,
but neither can it be resolved through
an assumption of unity.

A sequence of conditions may be associated
with one identity, multiple identities, or
an ontological configuration whose numerical
structure has not yet been established.

The fact that these conditions belong to
one continuous process does not prove
that they constitute one identity.

Similarly, the presence of multiple
manifestations, components, or relations
does not prove that multiple identities
are involved.

Unity of identity must be distinguished
from continuity of a process, unity of
a source, unity of an aggregate, and
continuity of a configuration.

The question of whether continuity requires
a numerically unified identity, and whether
a continuous ontological configuration
may involve more than one identity,
must remain open to investigation.

\subsection{The Possibility of Continuity Without Complete Preservation}

Identity continuity raises the question
of whether an identity can remain continuous
without the preservation of every condition
associated with it.

A strict preservation requirement would
risk identifying identity continuity with
the unchanged persistence of all properties,
determinations, relations, or components.

Such a requirement cannot be adopted
without establishing that every relevant
feature is constitutive of identity.

It is possible that some conditions may
change while others remain relevant to
the preservation of identity. It is also
possible that the significance of a change
depends upon the ontological domain or
respect in which it occurs.

However, this does not establish that
identity continuity is compatible with
every possible transformation.

The conditions under which preservation
can occur despite change must therefore
be investigated independently.

\subsection{The Possibility of Discontinuity Without Immediate Non-Identity}

The converse problem is equally important.

An interruption in a particular process,
manifestation, or configuration does not
automatically establish the numerical
non-identity of the conditions preceding
and following that interruption.

A discontinuity may concern one ontological
relation without necessarily determining
the identity relation.

For example, the interruption of a
manifestation or the cessation of a
particular process may leave unresolved
whether the identity associated with
the relevant existent persists.

This possibility must not be confused
with the claim that identity necessarily
survives every interruption.

Rather, it establishes that discontinuity
in one respect cannot be generalized
into discontinuity of identity without
an additional argument.

The subject and respect of discontinuity
must therefore be specified before
conclusions about identity are drawn.

\subsection{The Independence of Identity Continuity from Recognition}

The present inquiry concerns ontological
continuity rather than the recognition
of continuity by a knowing subject.

Identity continuity must not be defined
through memory, recollection, recognition,
observation, or the ability to identify
successive conditions as belonging to
one existent.

Such activities may provide epistemic
grounds for making claims about continuity,
but they do not constitute the ontological
conditions being investigated here.

An identity may be continuous or
discontinuous independently of whether
its continuity is recognized.

Likewise, the failure to identify a
continuous relation does not establish
that no such relation exists.

The distinction between ontological
continuity and epistemic access to
continuity is therefore essential.

\subsection{The Subject and Respect of Identity Continuity}

Every claim concerning identity continuity
must specify its subject and respect.

The subject may be an existent, an identity,
a configuration, a determination, or
another ontological aspect.

The respect may concern temporal succession,
structural transformation, material
constitution, relational organization,
manifestation, or another specified
condition.

Without this specification, statements
such as ``the existent continues,''
``the identity persists,'' or
``the configuration remains continuous''
may appear equivalent while referring
to different ontological questions.

A rigorous analysis must therefore
identify what is claimed to continue,
what changes, and which relation
connects the relevant conditions.

This requirement also prevents conclusions
about one ontological aspect from being
unjustifiably transferred to the identity
of the existent as a whole.

\subsection{Methodological Requirements for the Inquiry}

The investigation of identity continuity
must satisfy several methodological
requirements:

\begin{enumerate}
    \item Identity must be distinguished
    from persistence and continuity.

    \item Temporal succession must not
    be treated as sufficient evidence
    of identity continuity.

    \item Qualitative difference must
    be distinguished from numerical
    non-identity.

    \item Invariance must not be assumed
    to be either a sufficient or a
    necessary condition of identity
    continuity without argument.

    \item The continuity of manifestation
    must be distinguished from the
    continuity of identity.

    \item Neither unity nor plurality
    of identity may be presupposed.

    \item The interruption of a process
    or manifestation must not be
    generalized automatically into
    the interruption of identity.

    \item The absence of an identified
    continuity relation must not be
    confused with proof of ontological
    discontinuity.

    \item Every claim must specify
    its subject, domain, and respect.

    \item The analysis must remain
    independent of observation,
    memory, recognition, and
    epistemic accessibility.
\end{enumerate}

These requirements establish the
conceptual boundaries within which
the problem of identity continuity
can be investigated.

\subsection{Scope of the Present Section}

The present section has formulated
the ontological problem of identity
continuity and established its
distinction from several related
concepts.

It has not yet provided a complete
account of persistence, invariance,
change, or the conditions under
which identity may be preserved.

Nor has it established that every
existent possesses a single identity
that necessarily continues through
all transformations.

The subsequent sections will examine
these problems in greater detail.

The next section focuses on the
relation between identity and
persistence, clarifying whether
the continued existence of an
existent can provide sufficient
grounds for establishing the
continuity of its identity.

% =============================================================================
% End of Section 01
% =============================================================================

% ============================================================
% Section 02 — Identity and Persistence
% Book I: Ontology of Viability
% Chapter 05: Identity and Continuity
%
% Purpose:
%   To distinguish the persistence of an existent from
%   the persistence of its identity and from the continuity
%   of identity across temporal succession.
%
% Conceptual Policy:
%   1. Identity is not reducible to persistence.
%   2. Persistence of an existent does not automatically
%      establish persistence of identity.
%   3. Persistence, continuity, and temporal succession
%      must remain analytically distinct.
%   4. Material, structural, functional, relational, and
%      manifestational preservation are not independently
%      sufficient criteria of identity persistence.
%   5. Neither the termination nor the preservation of
%      identity may be inferred solely from the interruption
%      or continuation of a particular manifestation.
% ============================================================

\section{Identity and Persistence}
\label{sec:identity-and-persistence}

\subsection{The Ontological Distinction Between Identity and Persistence}

Identity and persistence are closely related ontological
notions, but they do not designate the same condition.

Identity concerns the numerical and ontological identity
of an existent: the question of what makes an existent
the existent that it is, without reducing this question
to any particular determination, configuration, or
temporal condition.

Persistence concerns the continuation of an existent
through a specified interval or under specified conditions.
It addresses whether an existent remains in existence
across a domain of temporal or other ordered variation.

The distinction is necessary because the continuation
of an existent and the persistence of its identity are
not automatically equivalent.

An existent may persist through a temporal interval while
undergoing substantial alterations in its determinations,
structure, relations, material constitution, or
manifestations. Conversely, the cessation of a particular
configuration or manifestation does not, by itself,
establish that the identity associated with it has ceased
to exist.

Persistence therefore cannot serve as an immediate
substitute for an ontological account of identity.

The present section examines the relation between these
notions without assuming that persistence necessarily
preserves identity or that the interruption of persistence
necessarily establishes the termination of identity.

\subsection{Persistence of an Existent and Persistence of Identity}

A distinction must be established between the persistence
of an existent and the persistence of its identity.

The persistence of an existent concerns the continuation
of that existent through a specified interval or domain.
The persistence of identity concerns whether the numerical
identity of that existent remains the same across the
relevant transformations or conditions.

These questions are related, but the answer to one does
not automatically determine the answer to the other.

The persistence of an existent may be established in
a particular respect without resolving whether its
identity has remained numerically the same. For example,
the continued presence of a material configuration,
the uninterrupted operation of a process, or the
maintenance of a recognizable structural organization
may establish certain forms of persistence. None of
these facts alone establishes the persistence of identity.

Likewise, a failure to establish the uninterrupted
persistence of a particular manifestation does not
necessarily establish the non-persistence of identity.

The distinction prevents the ontological question from
being settled through the mere observation of temporal
continuation or interruption.

Accordingly, every claim concerning persistence must
specify what is said to persist: the existent, its identity,
a determination, a structure, a relation, a process,
or a manifestation.

\subsection{Persistence and Temporal Continuity}

Persistence and temporal continuity must also be
distinguished.

Persistence refers to the continuation of an existent
through a specified interval. Temporal continuity concerns
the relation between successive temporal conditions
and whether the relevant succession is uninterrupted
or otherwise continuous according to a defined criterion.

An existent may be described as persistent even when
a particular process or manifestation associated with
it is interrupted. Similarly, temporal continuity in
a process does not necessarily establish the persistence
of the identity of the existent involved in that process.

The continuity of a process and the persistence of
an existent therefore belong to different analytical
levels.

A process may continue while its underlying conditions
or participating entities change. An existent may also
remain in existence despite interruptions in particular
activities or manifestations.

It follows that uninterrupted temporal succession
cannot be treated as a universal criterion of identity
persistence.

Nor can the presence of an interruption be treated
as sufficient evidence that identity has ceased.

The ontological significance of an interruption depends
on what has been interrupted, the domain in which the
interruption occurs, and the respect under which
persistence is being examined.

\subsection{Persistence and the Continuity of Manifestation}

The persistence of identity must not be conflated with
the continuity of manifestation.

Manifestations may continue, change, diminish, become
interrupted, or cease under different conditions.
Their continuity concerns the persistence or succession
of particular forms through which an existent is manifested.

Identity, however, cannot be identified with the
continuous presence of any particular manifestation.

The continued existence of a manifestation does not
independently establish that the identity associated
with it has persisted. A manifestation may remain
continuous while the ontological identity of its source
or bearer is transformed or otherwise brought into question.

Conversely, the interruption or absence of a manifestation
does not independently establish that the corresponding
identity has ceased to persist.

This distinction also applies to the capacity to manifest.
The continuity of the capacity to manifest is analytically
different from the continuity of actual manifestations
and from the persistence of identity.

An existent may retain a capacity without continuously
actualizing it. A manifestation may be interrupted
without thereby establishing the disappearance of
the capacity associated with it. Yet the persistence
of such a capacity, considered alone, does not establish
the persistence of identity.

Three questions must therefore remain distinct:

\begin{enumerate}
    \item Does the identity of the existent persist?
    \item Do its manifestations continue?
    \item Does its capacity to manifest continue?
\end{enumerate}

No one of these questions can be answered solely by
answering either of the other two.

\subsection{Persistence Through Material and Structural Change}

The persistence of an existent may be examined in relation
to changes in its material constitution and structural
organization.

Material continuity is one possible respect under which
persistence may be investigated. An existent may retain
some of its material constituents across a specified
interval, undergo partial material replacement, or
experience extensive material transformation.

However, the persistence of material constituents does
not by itself establish the persistence of identity.
Nor does the replacement or alteration of material
constituents automatically establish the emergence
of a numerically different identity.

The same caution applies to structural persistence.

A structure may remain relatively stable while its
components, relations, or material realization change.
Alternatively, a structure may undergo substantial
reorganization while the question of the identity of
the existent remains unresolved.

Structural preservation and structural transformation
are therefore relevant to the investigation of persistence,
but neither can be treated as a universal and sufficient
criterion of identity.

The ontological significance of material and structural
change must be determined in relation to the specific
existent and the relevant conditions of identity.

\subsection{Persistence Through Functional and Relational Change}

Persistence may also be investigated through functional
and relational conditions.

An existent may retain a function while undergoing
changes in its internal organization or external relations.
It may also acquire new functions, lose previous functions,
or enter into different relational configurations.

Functional continuity does not independently establish
identity persistence. The same function may be realized
by numerically different existents, and a single existent
may undergo functional transformation without this fact
alone establishing a change of identity.

Likewise, relational continuity cannot be treated as
a sufficient criterion of identity persistence.

An existent may remain related to the same entities
while undergoing transformations in other respects.
It may also enter into new relations or lose existing
relations without the ontological question of identity
being thereby settled.

Functions and relations may contribute to the conditions
under which persistence is investigated. Their significance,
however, must be established rather than assumed.

The persistence of a role, function, or relational position
must therefore be distinguished from the persistence
of the identity of the existent occupying that role
or position.

\subsection{Interruption, Absence, and the Problem of Persistence}

The problem of persistence becomes particularly difficult
when an existent or one of its manifestations undergoes
interruption.

An interruption may concern a process, an activity,
a manifestation, a relation, or a particular condition
of existence. These forms of interruption must not
be treated as ontologically equivalent.

The cessation of a process does not necessarily establish
the cessation of the existent associated with it.
The absence of a manifestation does not necessarily
establish the absence of its source or bearer.
Likewise, the discontinuation of a particular relation
does not necessarily establish the termination of
the identity of the entities involved.

The possibility of interruption therefore requires
a distinction between the interruption of a particular
mode of persistence and the termination of identity.

A further question concerns whether an existent can
persist through an interval during which some of its
manifestations or processes are absent.

This possibility cannot be rejected merely because
the relevant manifestation is not continuously present.
Nor can it be affirmed without examining the ontological
conditions of the existent concerned.

The relation between interruption and persistence
must therefore remain subject to the domain and respect
under examination.

\subsection{Temporary Cessation and the Possibility of Resumption}

The possibility of resumption introduces a further
distinction between uninterrupted persistence and
persistence across an interruption.

A process or manifestation may cease and subsequently
resume. The resumption may involve the continuation
of a previously existing process, the emergence of
a new process, or a relation between distinct processes.
The mere similarity between the conditions before
and after the interruption does not settle which of
these possibilities obtains.

Likewise, the resumption of a manifestation does not
by itself establish that the identity associated with
the earlier manifestation has persisted throughout
the interval.

Nor does the interruption necessarily establish that
the identity has ceased and subsequently been recreated.

The ontological relation between the conditions before
and after an interruption requires independent examination.
It cannot be resolved solely through the temporal
distance between them, the similarity of their
manifestations, or the restoration of a previous
configuration.

The possibility of persistence across interruption
must therefore be distinguished from the uninterrupted
continuity of a process or manifestation.

\subsection{Persistence and Causal Continuity}

Persistence must also be distinguished from causal
continuity.

A causal sequence may connect successive conditions,
events, or configurations. Such a sequence may explain
how one condition produces or contributes to another.
However, causal connection alone does not establish
that the numerically same identity persists across
the sequence.

Causal continuity may connect distinct existents,
and the persistence of an existent cannot be inferred
solely from the existence of a causal relation between
successive conditions.

Conversely, the absence of an established causal
sequence does not, without further argument, settle
the question of identity persistence.

Causal relations may provide relevant conditions
for investigating persistence, but they cannot replace
the ontological analysis of identity.

It is therefore necessary to distinguish the continuity
of causal production from the persistence of the identity
of the entities or processes involved.

\subsection{Persistence and the Conditions of Identity}

The persistence of identity cannot be established
through the isolated preservation of a single feature
without an adequate account of its ontological relevance.

Material constitution, structural organization,
functional capacity, relational configuration,
manifestational continuity, and causal connection
may all be relevant to the investigation. Nevertheless,
none of them can be assumed in advance to constitute
a universally sufficient condition of identity persistence.

The same feature may have different significance
in different domains or for different kinds of existents.
A condition that is essential to the persistence of
one existent may be non-essential to another.

Furthermore, the preservation of a feature in one
respect does not establish the preservation of identity
in every respect.

The investigation must therefore distinguish between:

\begin{itemize}
    \item conditions that accompany persistence;
    \item conditions that contribute to persistence;
    \item conditions that may be necessary for persistence
    in a specified domain;
    \item conditions that may be sufficient for persistence
    under explicitly defined circumstances;
    \item conditions that are merely correlated with
    the persistence of identity.
\end{itemize}

These categories cannot be treated as interchangeable.

The ontological status of each proposed condition must
be established through argument. No criterion of identity
persistence should be adopted merely because it describes
a frequently observed form of continuity.

\subsection{The Non-Equivalence of Persistence and Identity Preservation}

The preceding distinctions establish that persistence
and identity preservation are related but non-equivalent
notions.

Persistence may describe the continuation of an existent,
a process, a manifestation, a structure, or a relation.
Identity preservation concerns the continued numerical
identity of the existent under examination.

The persistence of one of these conditions does not
automatically entail the persistence of the others.

This non-equivalence has two important consequences.

First, the continued existence of something through
a specified interval cannot, without further argument,
be treated as proof that its identity has remained
numerically the same.

Second, the interruption or cessation of a particular
form of persistence cannot, without further argument,
be treated as proof that identity has ceased.

The relation between persistence and identity preservation
must therefore be established through an account of
the relevant ontological conditions, rather than through
the unqualified use of the term persistence.

\subsection{The Subject and Respect of Persistence}

Every ontological claim concerning persistence must
identify its subject and its respect.

The subject may be an existent, an identity, a process,
a structure, a material configuration, a relation,
a capacity, or a manifestation.

The respect may concern temporal duration, material
constitution, structural organization, functional
continuity, relational configuration, causal succession,
or another explicitly defined domain.

Without these distinctions, claims of persistence
remain ambiguous.

For example, the statement that an existent persists
may refer to its continued existence, the preservation
of its numerical identity, the continuity of its
manifestations, or the uninterrupted operation of
a process. These claims are not equivalent and require
different forms of ontological justification.

The subject and respect of persistence must therefore
be specified before any conclusion concerning identity
preservation can be drawn.

\subsection{Methodological Requirements}

The investigation of identity and persistence must
observe several methodological requirements.

First, persistence must not be used as a substitute
definition of identity.

Second, the persistence of an existent must be distinguished
from the persistence of its identity.

Third, temporal continuity, causal continuity, and
continuity of manifestation must be examined separately.

Fourth, material, structural, functional, and relational
preservation must not be treated as universally sufficient
criteria of identity persistence.

Fifth, interruption, absence, and resumption must be
analyzed according to the particular domain and respect
involved.

Sixth, neither the continuation nor the termination
of identity may be inferred solely from the continuation
or cessation of a particular process or manifestation.

Finally, the ontological conditions of identity persistence
must remain open to investigation rather than being
determined in advance by a preferred model of continuity.

\subsection{Scope of the Present Section}

This section has distinguished identity from persistence
and examined the different senses in which an existent,
its identity, its manifestations, and its associated
processes may continue or cease.

The central conclusion is that persistence does not
provide an independent or sufficient explanation of
identity preservation.

The persistence of an existent, the persistence of
its identity, the continuity of its manifestations,
and the continuity of its capacity to manifest are
distinct ontological questions. Their relations must
be established through further analysis.

The next section examines the relation between invariance
and continuity. It investigates whether the preservation
of invariant features can contribute to an account
of identity continuity, and under what conditions
such preservation may or may not have ontological
significance.

% ============================================================
% End of Section 02
% ============================================================

% ============================================================
% Section 03 — Invariance and Continuity
% Book I: Ontology of Viability
% Chapter 05: Identity and Continuity
%
% Purpose:
%   To examine the ontological relation between invariance
%   and continuity without reducing identity continuity
%   to the preservation of invariant features.
%
% Conceptual Policy:
%   1. Invariance is not identity.
%   2. Invariance is not persistence or continuity.
%   3. The preservation of an invariant does not by itself
%      establish the continuity of identity.
%   4. The absence of an identified invariant does not
%      by itself establish discontinuity of identity.
%   5. Every invariant must be specified relative to a
%      transformation, domain, and respect.
%   6. Formal invariance must not automatically be treated
%      as ontological invariance.
% ============================================================

\section{Invariance and Continuity}
\label{sec:ch05-invariance-and-continuity}

\subsection{The Ontological Problem of Invariance and Continuity}

The relation between invariance and continuity constitutes
a central problem in the investigation of identity
through change.

An existent may undergo transformations while preserving
certain features, relations, structures, or determinations.
Such preservation may be described in terms of invariance.
The existence of invariants can therefore provide relevant
information about the conditions under which an existent
continues through transformation.

However, the preservation of an invariant and the
continuity of identity are not equivalent ontological
conditions.

An invariant may remain preserved across a specified
transformation without establishing that the numerically
same identity has continued. Conversely, an existent
may undergo transformations in which no previously
identified invariant remains preserved, without this
fact alone establishing that its identity has ceased.

The problem is therefore not simply whether something
remains invariant, but whether the invariant in question
has a justified ontological relation to the identity
whose continuity is being examined.

This requires distinguishing the formal preservation
of a feature from the ontological preservation of
identity. It also requires identifying the transformations
under consideration and the domain and respect in which
invariance is asserted.

The present section examines these relations without
assuming that identity is constituted by invariance
or that continuity can be established through the
preservation of invariant features alone.

\subsection{Invariance as a Relation to Transformation}

Invariance is meaningful only in relation to a specified
class of transformations.

A feature, relation, structure, or determination is
invariant when it remains preserved under the relevant
transformations according to an explicitly defined
criterion.

Consequently, an assertion of invariance requires
at least three elements:

\begin{itemize}
    \item a feature or property whose preservation
    is being examined;
    \item a transformation or class of transformations
    relative to which preservation is assessed;
    \item a domain and respect in which the invariance
    is asserted.
\end{itemize}

Without these specifications, the claim that something
is invariant remains incomplete.

A feature may remain invariant under one class of
transformations while changing under another. A relation
may be preserved in one domain while failing to remain
preserved in a different domain. A structure may remain
formally equivalent under a transformation even though
its material realization changes.

Invariance is therefore not an absolute description
of an existent independently of the transformations
and respects under consideration.

This relativity to transformation is particularly
important when invariance is invoked in discussions
of identity continuity. The preservation of a feature
under a limited class of transformations cannot be
extended without justification to all possible
transformations of the existent.

\subsection{The Distinction Between Invariance and Continuity}

Invariance and continuity address different questions.

Invariance concerns the preservation of a specified
feature under a specified transformation.

Continuity concerns whether a relevant existent,
identity, process, or manifestation continues across
a specified succession or domain of change.

The first is therefore a relation between a feature
and a transformation. The second concerns the continuation
of a subject through a relevant domain.

An invariant may be preserved across a transformation
without establishing the continuity of the existent
to which it is attributed.

Similarly, continuity may obtain in a relevant sense
even when some features undergo alteration and no
single previously identified invariant remains preserved.

The two notions may be related, but their relation
must be established rather than assumed.

The preservation of an invariant can contribute to
an account of continuity when the invariant has been
shown to be ontologically relevant to the subject
under examination. Yet the mere identification of
a preserved feature does not establish this relevance.

Likewise, the transformation of a feature does not
automatically establish discontinuity. It may indicate
a change in a particular determination or respect
without determining whether the identity of the
existent has continued.

\subsection{Invariance and the Persistence of Identity}

The persistence of identity must not be defined as
the preservation of an invariant.

Such a definition would make the ontological question
of identity dependent upon a criterion that has not
yet been independently justified.

If identity were defined through the preservation
of a selected invariant, the subsequent claim that
identity persists because the invariant remains
preserved would merely repeat the initial definition.
It would not establish an independent ontological
account of identity continuity.

Furthermore, the selection of an invariant requires
an explanation of why that particular feature should
be regarded as relevant to identity.

A material feature, structural relation, functional
organization, or relational configuration may remain
preserved across transformations. However, the preservation
of any such feature does not automatically establish
that the identity associated with it has persisted.

The same invariant may characterize distinct existents.
Consequently, the preservation of a feature cannot
by itself establish numerical identity between the
subjects that exhibit it.

Conversely, the transformation or disappearance of
a particular invariant does not independently establish
that the identity of its bearer has ceased.

The relation between invariance and identity persistence
must therefore be investigated without making either
notion a substitute for the other.

\subsection{The Problem of Selecting Identity-Relevant Invariants}

Not every invariant has the same ontological significance.

A feature may remain preserved because it belongs
to a relatively stable configuration, because the
considered transformation does not affect it, or
because the criterion of invariance is defined at
a level of abstraction that excludes certain changes.

None of these circumstances independently establishes
that the feature is constitutive of identity.

The identification of an identity-relevant invariant
therefore requires more than demonstrating preservation.

It requires an account of the relation between the
invariant and the existent whose identity is under
examination.

At least three questions must be distinguished:

\begin{enumerate}
    \item Is the feature genuinely invariant under
    the transformations being considered?
    \item Is the feature ontologically relevant to
    the identity of the existent?
    \item Does the preservation of this feature provide
    sufficient grounds for establishing continuity
    of identity?
\end{enumerate}

An affirmative answer to the first question does
not entail affirmative answers to the second or third.

A feature may be invariant without being identity-relevant.
It may be identity-relevant without being sufficient
to establish identity continuity independently.
It may also contribute to an account of identity
preservation only in conjunction with other conditions.

The ontological role of an invariant must therefore
be established for the particular existent, domain,
and class of transformations under examination.

\subsection{The Non-Equivalence of Formal and Ontological Invariance}

Formal invariance and ontological invariance must
remain distinct.

A formal system may identify quantities, relations,
or structures that remain unchanged under a specified
transformation. Such invariance is established relative
to the definitions, rules, and domain of the formal
system.

Ontological invariance concerns the preservation of
a feature or determination in the reality of an existent
under the relevant transformations.

Formal invariance may provide a useful means of
representing or investigating ontological relations.
Nevertheless, the preservation of a mathematical
or formal structure does not automatically establish
that the corresponding ontological feature has been
preserved.

A formal mapping may preserve a relation while the
entities represented by that relation undergo changes
whose ontological significance is not captured by
the mapping.

Likewise, formal equivalence between two configurations
does not independently establish that they belong
to one and the same numerical identity.

The transition from formal invariance to ontological
invariance therefore requires an argument establishing
the adequacy of the formal representation for the
ontological subject under examination.

Without such an argument, formal preservation remains
a result internal to the formal domain.

\subsection{Transformation Classes and the Scope of Continuity}

The relation between invariance and continuity depends
upon the class of transformations being considered.

An existent may preserve a feature under minor structural
alterations while failing to preserve it under more
extensive transformations. A relation may remain
invariant under transformations of material realization
but not under transformations of relational organization.

The continuity of identity cannot therefore be inferred
from invariance under a restricted class of transformations
unless the relevance and adequacy of that class have
been established.

The distinction also prevents the unrestricted use
of the proposition that identity remains invariant
under all temporal transformations.

Temporal transformations do not constitute a
self-justifying class of transformations under which
identity must necessarily remain invariant.

The fact that a transformation occurs through time
does not establish that it is compatible with the
persistence of identity. Nor does it establish that
it is incompatible with such persistence.

The ontological status of each transformation must
be examined independently.

A claim of continuity through transformation must
therefore specify which transformations are involved,
which features remain preserved, and why those features
are relevant to the identity under examination.

\subsection{Continuity Without an Identified Invariant}

The absence of an identified invariant does not
necessarily establish the discontinuity of identity.

There may be several reasons why an invariant has
not been identified.

The relevant feature may not yet be known. The
transformation class may not have been adequately
specified. The chosen level of analysis may fail
to capture the relevant ontological relation.
Alternatively, identity continuity may not depend
upon the preservation of a single invariant feature.

Consequently, the failure to identify an invariant
must be distinguished from the demonstrated absence
of every possible identity-relevant invariant.

Even the demonstrated transformation of all previously
identified features does not automatically establish
that the numerical identity of the existent has ceased.

This does not imply that identity must persist whenever
no invariant is identified. It means only that the
absence of an identified invariant is insufficient
as an independent proof of discontinuity.

The question remains open until the relevant ontological
conditions have been examined.

\subsection{Invariance Without Continuity of Identity}

The reverse possibility must also be considered.

An invariant may remain preserved across a transformation
even when the numerical identity of the subject
has not been established as continuous.

Distinct existents may share the same invariant feature.
A structure may be reproduced in another existent.
A relation may be maintained across a transition
between different subjects.

In these circumstances, the preservation of the
invariant establishes a form of continuity or
equivalence at the level of the relevant feature,
but it does not necessarily establish continuity
of numerical identity.

This possibility demonstrates why invariance cannot
serve as a sufficient criterion of identity persistence.

The ontological question concerns not merely whether
a feature remains the same, but whether the existent
whose identity is being examined remains numerically
the same existent.

The distinction between feature preservation and
identity preservation must therefore be maintained
even when the invariant is highly stable or structurally
significant.

\subsection{Invariance, Change, and the Limits of Preservation}

Invariance must not be understood as the absence
of all change.

An existent may preserve certain features while
undergoing transformations in others. The preservation
of one respect does not imply the immutability of
the whole existent.

Likewise, the alteration of a feature does not
necessarily establish that every relevant form of
continuity has been interrupted.

The ontological significance of change depends upon
the subject and respect under examination.

A transformation may preserve an invariant at one
level while altering a different determination at
another. It may preserve a structural relation while
changing its material realization. It may preserve
a relational organization while altering the functions
performed within it.

These possibilities demonstrate that invariance and
change are not mutually exclusive descriptions
of an existent considered without qualification.

The preservation of selected features must not be
generalized into a claim of total preservation.
Likewise, the alteration of selected features must
not be generalized into a claim of total discontinuity.

An adequate account of identity continuity must
therefore specify the respects in which preservation
and transformation occur.

\subsection{The Relation Between Invariance and Manifestation Continuity}

The continuity of manifestations may involve the
preservation of certain invariant features.

A sequence of manifestations may preserve a structural
relation, a material characteristic, or another
specified determination across successive conditions.

Such preservation may help describe the continuity
of the manifestation sequence. Nevertheless, it
does not independently establish the persistence
of the identity associated with those manifestations.

The same invariant feature may occur in manifestations
belonging to distinct existents. Conversely, manifestations
associated with one existent may undergo transformations
that alter previously preserved features.

The relation between manifestation continuity and
identity continuity must therefore be established
independently of the mere preservation of manifest
features.

The continuity of a manifestation sequence may be
relevant to the investigation of identity, but it
cannot be treated as a direct substitute for it.

\subsection{The Subject and Respect of Invariance}

Every claim concerning invariance and continuity
must identify its subject and respect.

The subject may be an existent, an identity, a
determination, a structure, a relation, a material
configuration, or a manifestation.

The respect must specify the feature whose preservation
is being examined and the transformations under
which that preservation is asserted.

An invariant identified at the level of structure
cannot automatically be transferred to the level
of identity. Nor can the preservation of a feature
in one domain be generalized to every domain in
which the existent undergoes transformation.

The same discipline applies to continuity. A claim
that continuity has been established must identify
what continues and in which respect.

Without these specifications, the relation between
invariance and continuity remains ambiguous and
cannot support a determinate ontological conclusion.

\subsection{Methodological Requirements}

The investigation of invariance and continuity must
observe the following methodological requirements:

\begin{enumerate}
    \item Invariance must be specified relative to
    an identifiable transformation or class of
    transformations.
    \item The domain and respect of every invariance
    claim must be explicitly identified.
    \item Formal invariance must not automatically
    be treated as ontological invariance.
    \item The preservation of an invariant must not
    be treated as sufficient proof of identity continuity.
    \item The absence of an identified invariant must
    not be treated as sufficient proof of identity
    discontinuity.
    \item The preservation of a feature must be
    distinguished from the preservation of numerical
    identity.
    \item The relation between invariance and identity
    must be established independently rather than
    assumed through definition.
\end{enumerate}

These requirements prevent invariance from becoming
an unrestricted criterion through which every question
of identity continuity is prematurely resolved.

\subsection{Scope of the Present Section}

This section has examined the relation between
invariance and continuity while preserving their
ontological distinction.

Invariance concerns the preservation of a specified
feature under a specified class of transformations.
Continuity concerns the continuation of a subject
across a relevant domain or succession. Neither
notion can be reduced to the other.

The preservation of an invariant does not independently
establish continuity of identity, and the absence
of an identified invariant does not independently
establish discontinuity of identity.

The ontological relevance of invariants must be
demonstrated in relation to the identity under
examination, the transformations involved, and
the domain and respect of the claim.

The next section examines identity and non-identity
through difference. It investigates whether qualitative,
structural, material, or relational differences
between successive conditions can establish numerical
non-identity, and why difference must not be treated
as an immediate proof of the termination or replacement
of identity.

% ============================================================
% End of Section 03
% ============================================================

% ============================================================
% Section 04 — Identity and Non-Identity Through Difference
% Book I: Ontology of Viability
% Chapter 05: Identity and Continuity
%
% Purpose:
%   To examine whether difference between successive
%   conditions can establish numerical non-identity,
%   and to distinguish qualitative difference from
%   ontological discontinuity.
%
% Conceptual Policy:
%   1. Difference is not non-identity.
%   2. Qualitative transformation does not automatically
%      establish numerical replacement.
%   3. Similarity does not independently establish identity.
%   4. Difference must be specified by subject, domain,
%      and respect.
%   5. Neither identity nor non-identity should be inferred
%      from a single unqualified comparison.
% ============================================================

\section{Identity and Non-Identity Through Difference}
\label{sec:identity-and-non-identity-through-difference}

\subsection{The Ontological Problem of Difference and Non-Identity}

The relation between difference and non-identity constitutes
a fundamental problem in the investigation of identity
through transformation.

An existent may undergo qualitative, material, structural,
functional, relational, or manifestational differences
across successive conditions. Such differences may be
substantial and may alter the way in which the existent
is determined or manifested.

Nevertheless, the presence of difference does not
immediately establish numerical non-identity.

The proposition that an existent differs from a previous
condition and the proposition that it is no longer
the same existent are distinct ontological claims.
The first concerns a difference between determinations
or conditions. The second concerns the numerical identity
of the existent across the relevant succession.

Confusing these claims would make every transformation
a possible termination of identity merely because
some feature has changed. Such a conclusion would
require an independent argument establishing that
the relevant difference is incompatible with the
persistence of identity.

The opposite inference is equally problematic.
The absence of an identified difference, or the
presence of extensive similarity, does not independently
establish numerical identity.

The ontological problem is therefore not whether
difference exists, but whether a particular difference
has the significance required to establish non-identity.

\subsection{Qualitative Difference and Numerical Non-Identity}

Qualitative difference concerns the distinction between
features, determinations, properties, or conditions.

Numerical non-identity concerns whether two subjects
are distinct existents rather than one and the same
existent considered across different conditions.

These notions belong to different analytical levels.

Two conditions of an existent may differ qualitatively
while remaining conditions of one and the same existent.
A material configuration may change, a structural
organization may be modified, or a functional capacity
may be transformed without these changes alone
establishing numerical replacement.

Likewise, two numerically distinct existents may share
many or even all of the features selected for a particular
comparison.

Qualitative similarity and difference therefore cannot
be treated as direct measures of numerical identity.

The distinction requires that every comparison specify
whether it concerns features, configurations, relations,
manifestations, or the numerical identity of the
subjects being compared.

Without this distinction, qualitative descriptions
may be incorrectly transformed into ontological
conclusions.

\subsection{Difference Across Temporal Conditions}

The problem becomes particularly significant when
difference is examined across temporal conditions.

An existent may be described at different temporal
stages through distinct configurations or determinations.
The earlier and later conditions may differ in numerous
respects while still raising the question of whether
they belong to one and the same identity.

The mere fact that the conditions occur at different
times does not establish that they belong to different
existents. Nor does temporal succession establish
that they belong to one numerically identical existent.

The relation between successive conditions must be
examined rather than assumed.

A temporal sequence may involve transformation within
one existent, the succession of distinct existents,
or a more complex relation in which the relevant
subjects and respects have not yet been adequately
determined.

Consequently, the comparison of temporal stages
cannot be resolved through difference alone.

The temporal distance between two conditions, their
degree of similarity, or the magnitude of their
qualitative transformation does not independently
settle the question of numerical identity.

\subsection{Difference in Material Constitution}

Material difference is one important form of difference
through which identity continuity is often examined.

An existent may undergo material alteration through
the loss, replacement, addition, or reorganization
of material constituents.

Such changes may be extensive. Nevertheless, the
degree of material difference does not by itself
determine whether the identity of the existent has
persisted.

The preservation of material constituents does not
independently establish numerical identity. Different
existents may share the same material constituents
or material composition under a specified description.

Conversely, the replacement of some or many material
constituents does not automatically establish that
the existent has become numerically different.

The ontological significance of material transformation
depends upon the kind of existent, the relevant
conditions of its constitution, and the respect
under which identity is being examined.

Material difference must therefore be distinguished
from numerical non-identity.

\subsection{Difference in Structural Organization}

Structural difference concerns alterations in the
organization of components, relations, or internal
configurations.

An existent may undergo structural reorganization
while preserving some of its previous relations.
It may also acquire a substantially different
organization through transformation.

Neither structural preservation nor structural
difference independently determines numerical identity.

A similar structure may be realized by distinct
existents. A structure may also undergo transformation
without necessarily establishing the replacement
of the existent whose organization has changed.

The persistence of a structure and the persistence
of identity are therefore separate questions.

Structural difference may be relevant to an account
of identity when the structure has been shown to
play an ontologically significant role. Yet its
relevance cannot be assumed merely because structural
features are comparatively stable or readily described.

The transition from structural difference to
non-identity requires an account of why the particular
transformation is incompatible with the persistence
of the identity under examination.

\subsection{Difference in Functional Determination}

Functional difference concerns changes in the
operations, capacities, roles, or activities associated
with an existent.

An existent may acquire new functions, lose previous
functions, or undergo a transformation in the manner
in which its capacities are actualized.

Functional continuity does not independently establish
identity, since distinct existents may perform the
same function.

Likewise, functional discontinuity does not independently
establish non-identity. An existent may undergo
functional transformation without the numerical
identity question being settled by that transformation.

The same function may be realized through different
structures, materials, or relations. Conversely,
the same existent may participate in different
functional arrangements across successive conditions.

A functional description therefore cannot serve
as a universal criterion of numerical identity.

The ontological relevance of functional difference
must be established in relation to the existent
and the conditions under which its identity is
being examined.

\subsection{Difference in Relational Configuration}

An existent may undergo changes in its relational
configuration.

It may enter into new relations, lose previous
relations, or participate in a different organization
of relations while retaining other determinations.

Such relational transformations may significantly
alter its position within a wider configuration.
However, a change in relational position does not
automatically establish a change in numerical identity.

The same relational position may be occupied by
different existents. An existent may also occupy
different relational positions across successive
conditions.

Relational difference must therefore be distinguished
from numerical non-identity.

A relational configuration may contribute to the
conditions of identity preservation in a particular
domain, but its significance must be independently
established. It cannot be treated as a universal
criterion without examining the ontological status
of the relations involved.

\subsection{Difference in Manifestation}

Manifestations associated with an existent may
differ substantially across successive conditions.

An existent may manifest different determinations,
configurations, capacities, or relations. It may
also undergo changes in the extent, form, or organization
of its manifestations.

Such differences may establish that the manifestations
are not qualitatively identical. They do not
automatically establish that the identity associated
with them is numerically different.

The distinction between identity and manifestation
is therefore essential to the analysis of difference.

A manifestation cannot be treated as an exhaustive
representation of the identity of the existent.
Consequently, a difference between manifestations
cannot be transferred without argument into a
difference between identities.

Likewise, similar manifestations do not independently
establish that their sources or bearers possess
the same numerical identity.

The relation between manifestation difference and
identity difference must remain an independent
ontological question.

\subsection{The Difference Between Change and Replacement}

Change and replacement must not be treated as
equivalent.

Change concerns the alteration of a subject or
one of its determinations. Replacement concerns
the emergence or substitution of a numerically
different subject in relation to a previous one.

A change may be extensive without necessarily
constituting replacement. Conversely, replacement
may occur without requiring a substantial qualitative
difference between the previous and subsequent
subjects.

This distinction prevents the magnitude of change
from becoming an automatic criterion of non-identity.

The claim that an existent has changed significantly
does not, by itself, establish that another existent
has replaced it. Such a conclusion requires an
account of the ontological conditions under which
the transformation constitutes replacement rather
than alteration within the same identity.

Similarly, the restoration of a previous configuration
does not necessarily establish the restoration
of the previous numerical identity.

Change, replacement, and restoration must therefore
remain analytically distinct.

\subsection{The Problem of Gradual Difference}

Gradual transformation introduces a further difficulty.

An existent may undergo a sequence of relatively
small changes, each of which differs only slightly
from the preceding condition. Across an extended
succession, however, the final condition may differ
substantially from the initial one.

The gradual character of the transformation does
not independently establish continuity of identity.
Nor does the accumulation of qualitative differences
automatically establish numerical non-identity.

The problem cannot be resolved merely by examining
the magnitude of each individual change or the total
difference between the initial and final conditions.

A sequence of small transformations may involve
continuity of one existent, the succession of distinct
existents, or another ontological configuration
that requires further determination.

The question of identity must therefore be examined
across the succession as a whole and in relation
to the ontological conditions relevant to the
existent under consideration.

The absence of a single decisive transformation
does not establish identity continuity. Equally,
the accumulation of differences does not independently
establish identity discontinuity.

\subsection{Similarity and the Problem of False Identity}

The distinction between difference and non-identity
must be accompanied by a distinction between similarity
and identity.

Two existents may share a substantial number of
features, possess similar structures, perform
equivalent functions, or exhibit comparable
manifestations without being numerically identical.

Similarity concerns a relation of resemblance
under selected respects. Numerical identity concerns
whether the subjects are one and the same existent.

The two claims cannot be substituted for one another.

A sufficiently high degree of similarity does not
establish numerical identity. Nor does the reproduction
of a configuration or structure establish that
the reproduced existent is numerically identical
with the original.

This distinction is particularly important when
identity continuity is examined through preserved
features. A preserved feature may characterize
both successive conditions and distinct existents.

Consequently, the identification of similarity
or equivalence must not be treated as a complete
account of identity.

\subsection{Difference and the Problem of Multiple Identities}

Difference may also arise in relation to the possibility
of multiple identities.

An existent may exhibit several distinct configurations,
manifestations, or relational determinations. These
differences do not automatically establish that
the existent possesses multiple identities.

Conversely, the apparent unity of an existent does
not independently establish that only one identity
is involved.

The relation between difference and plurality of
identities must therefore be examined without
assuming that every distinct manifestation corresponds
to a distinct identity or that every manifestation
belongs to one identity.

The question of whether a transformation concerns
different determinations of one identity or the
relation between multiple identities cannot be
resolved through qualitative difference alone.

This issue requires an independent investigation
of the ontological conditions under which identity
is individuated and preserved.

The present section therefore leaves the number
of identities open rather than deriving it from
the multiplicity or similarity of observed conditions.

\subsection{Difference Without a Determinate Conclusion}

There are cases in which a difference can be established
without the available ontological grounds being
sufficient to determine whether identity has persisted.

A material, structural, functional, or relational
transformation may be clearly identifiable while
its significance for numerical identity remains
undetermined.

In such cases, the appropriate conclusion is not
necessarily identity or non-identity. It may be
that the relevant difference has been established
but that its ontological implications remain unresolved.

This distinction is methodologically important.

The inability to establish identity continuity
does not automatically establish non-identity.
Likewise, the inability to establish non-identity
does not automatically establish continuity.

An unresolved question must remain unresolved until
adequate grounds are available.

The distinction between established difference
and undetermined identity prevents the analysis
from producing conclusions that exceed the evidence
and ontological argument available.

\subsection{The Subject and Respect of Difference}

Every claim concerning difference must specify
its subject and respect.

The subject may be an existent, an identity,
a determination, a manifestation, a structure,
a relation, or a temporal condition.

The respect may concern material constitution,
structural organization, functional determination,
relational configuration, or another explicitly
defined domain.

A difference established between two manifestations
cannot automatically be transferred to the identities
associated with them. A difference in material
constitution cannot automatically be generalized
to every other respect of the existent.

Likewise, the comparison of two temporal conditions
must distinguish between differences within one
existent and differences between distinct existents.

Without these specifications, the concept of difference
remains too broad to support a determinate conclusion
about numerical identity.

\subsection{Methodological Requirements}

The investigation of identity and non-identity through
difference must observe the following requirements:

\begin{enumerate}
    \item Qualitative difference must be distinguished
    from numerical non-identity.
    \item Similarity must not be treated as sufficient
    proof of numerical identity.
    \item The magnitude of transformation must not
    be used as an automatic criterion of replacement.
    \item Material, structural, functional, relational,
    and manifestational differences must be examined
    separately.
    \item Change must be distinguished from replacement
    and restoration.
    \item Gradual transformation must not be assumed
    to guarantee identity continuity or discontinuity.
    \item Difference between manifestations must not
    be transferred automatically to difference between
    identities.
    \item The possibility of multiple identities must
    remain open to independent ontological investigation.
    \item Every conclusion must specify its subject,
    domain, and respect.
\end{enumerate}

These requirements prevent difference from becoming
an unrestricted criterion through which identity
is either affirmed or denied without adequate
ontological justification.

\subsection{Scope of the Present Section}

This section has established the distinction between
qualitative difference and numerical non-identity.

An existent may undergo substantial differences
in its material constitution, structural organization,
functional determination, relational configuration,
or manifestations without these differences alone
establishing that its identity has ceased.

Conversely, similarity between successive conditions
does not independently establish that they belong
to one and the same numerical identity.

Difference is therefore relevant to the investigation
of identity continuity, but it cannot serve as
an independent and sufficient criterion of either
identity preservation or non-identity.

The ontological significance of difference must be
determined in relation to the subject, domain, and
respect under examination.

The next section investigates the conditions of
identity preservation. It examines whether and under
what circumstances particular ontological conditions
can justify the claim that identity persists through
transformation, without reducing identity to any
single material, structural, relational, functional,
or invariant feature.

% ============================================================
% End of Section 04
% ============================================================

% ============================================================
% Section 05 — Conditions of Identity Preservation
% Book I: Ontology of Viability
% Chapter 05: Identity and Continuity
%
% Purpose:
%   To investigate the ontological conditions under which
%   the preservation of identity through transformation
%   may be established.
%
% Conceptual Policy:
%   1. Identity preservation must not be reduced to a
%      single material, structural, functional, relational,
%      or invariant condition.
%   2. No proposed condition is universally necessary
%      or sufficient without independent justification.
%   3. Identity preservation must be distinguished from
%      persistence, similarity, and continuity of manifestation.
%   4. Conditions of identity preservation must be examined
%      relative to the existent, domain, and respect involved.
%   5. The possibility of transformation without identity
%      loss must remain distinct from an unrestricted
%      assumption of identity persistence.
% ============================================================

\section{Conditions of Identity Preservation}
\label{sec:conditions-of-identity-preservation}

\subsection{The Ontological Problem of Identity Preservation}

The question of identity preservation concerns the
conditions under which an existent remains numerically
the same existent through transformation.

The preceding sections have established that identity
cannot be reduced to persistence, invariance, or the
continuity of particular manifestations. They have
also shown that qualitative difference does not
automatically establish numerical non-identity.

These distinctions, however, do not yet provide
an account of the conditions under which identity
is preserved.

The problem requires an examination of the ontological
conditions that may justify the attribution of one
and the same identity across different conditions
of an existent.

Such conditions cannot be selected merely because
they are observable, relatively stable, or frequently
associated with continuity. Their relevance must
be established in relation to the identity under
examination.

Material constitution, structural organization,
determination, relational configuration, invariance,
and manifestation may all contribute to an account
of identity preservation. Nevertheless, none can
be treated in advance as a universal and sufficient
criterion.

The present section therefore investigates identity
preservation as an ontological problem whose conditions
must be established through analysis rather than
assumed through a preferred model of continuity.

\subsection{The Difference Between Identity Preservation and Feature Preservation}

Identity preservation must be distinguished from
the preservation of particular features.

An existent may preserve a feature while undergoing
transformations in other respects. It may also lose
or transform a previously present feature without
this fact alone establishing the termination of
its identity.

Feature preservation concerns the continuation
of a determination, relation, structure, or other
specified characteristic. Identity preservation
concerns the continued numerical identity of the
existent itself.

The two may be related, but their relation requires
independent justification.

A feature may be shared by numerically distinct
existents. Consequently, its preservation cannot
by itself establish that the existent exhibiting
it remains numerically identical with an earlier
existent.

Likewise, the transformation of a feature does
not independently establish that the existent
has become numerically different.

The investigation must therefore distinguish between
features that remain preserved and the ontological
conditions under which the identity of their bearer
is preserved.

\subsection{Identity Preservation and the Ontological Unity of the Existent}

An account of identity preservation must examine
the relation between identity and the existent
whose identity is being investigated.

Identity is not an additional property that is
attached to an otherwise independently constituted
existent. Nor is the existent merely an aggregate
of features whose preservation automatically
establishes the preservation of identity.

The ontological question concerns the continued
identity of the existent through the transformations
that affect its determinations and conditions.

This requires distinguishing the unity of identity
from the unity of a configuration, structure,
material aggregate, or manifestation.

An existent may undergo changes in the organization
of its components or relations without these changes
alone determining whether its identity has been preserved.

Conversely, the apparent unity of a configuration
does not independently establish the numerical
identity of the existent constituted or manifested
through it.

The relation between identity and the unity of
the existent must therefore be examined without
reducing one to the other.

The conditions of identity preservation must account
for the existent as an ontological subject rather
than merely identifying a collection of features
that remain unchanged.

\subsection{Material Conditions of Identity Preservation}

Material constitution may be relevant to the
preservation of identity in particular domains.

An existent may depend upon a specific material
constitution for its continued existence or for
the realization of certain determinations. In such
cases, material conditions may contribute to the
conditions under which identity preservation is
examined.

However, the relevance of material constitution
does not establish that identity is universally
identical with material continuity.

Material constituents may be replaced, reorganized,
or transformed while the question of identity
remains open. Conversely, the preservation of
the same material constituents does not independently
establish numerical identity.

The material conditions of identity preservation
must therefore be examined in relation to the
kind of existent and the ontological role played
by its material constitution.

A material condition may be necessary in a particular
domain without being universally necessary. It
may contribute to identity preservation without
being sufficient to establish it.

The ontological status of material conditions must
be determined through a specific account of the
existent under examination.

\subsection{Structural Conditions of Identity Preservation}

Structural organization may also contribute to
the preservation of identity.

An existent may depend upon particular relations
among its components, determinations, or internal
configurations. The alteration of these relations
may affect the conditions under which its identity
is maintained.

Nevertheless, structural preservation cannot
be treated as a universal criterion of identity.

A structure may be reproduced in a numerically
distinct existent. Structural equivalence therefore
does not independently establish numerical identity.

Likewise, structural transformation does not
automatically establish the termination of identity.
An existent may undergo reorganization while the
ontological question of identity preservation
remains unresolved.

The relevant question is not merely whether a
structure has remained unchanged, but whether
the structure in question has an ontologically
constitutive role in the identity of the existent.

Such a role must be demonstrated for the particular
domain under examination.

\subsection{Determination and Identity Preservation}

Determinations may contribute to the conditions
under which identity preservation is investigated.

An existent is not without determination, and
transformations in its determinations may alter
the ways in which its existence is constituted
or manifested.

However, the preservation of a particular determination
does not independently establish identity preservation.
Nor does the alteration of a determination automatically
establish non-identity.

Determinations may differ in their ontological
significance. Some may be constitutive under
particular conditions, while others may be contingent
or transformable without determining the numerical
identity of the existent.

The investigation must therefore distinguish
between the preservation of determinations and
the preservation of identity.

It must also avoid assuming that every determination
has the same relation to identity or that identity
can be reconstructed as the totality of preserved
determinations.

The role of determination in identity preservation
requires an independent ontological account.

\subsection{Relational Conditions of Identity Preservation}

An existent may be situated within a configuration
of relations that contributes to the conditions
of its existence.

Some relations may be constitutive of the particular
existent under examination. Others may change
without necessarily affecting its numerical identity.

The persistence of a relational configuration
may therefore be relevant to identity preservation,
but it cannot be treated as a universal criterion.

The same relational organization may be instantiated
by distinct existents. Conversely, an existent
may undergo relational transformation without
necessarily becoming numerically different.

The ontological relevance of relations must be
established according to their specific role,
domain, and respect.

It is also necessary to distinguish the preservation
of a relation from the preservation of the identity
of the relata. A relation may persist even when
one or more of its relata change, and the continuity
of a relation does not independently establish
the numerical identity of the entities involved.

Relational conditions may contribute to an account
of identity preservation, but their significance
must not be assumed in advance.

\subsection{Invariance as a Possible Condition of Identity Preservation}

Invariance may provide a relevant condition in
the investigation of identity preservation.

A feature that remains invariant under a specified
class of transformations may indicate that certain
ontological conditions have been preserved.

However, invariance is not identical with identity
preservation.

An invariant may be shared by distinct existents,
and the preservation of an invariant does not
independently establish numerical identity.
Conversely, the absence of an identified invariant
does not necessarily establish that identity has
ceased to persist.

The relevance of an invariant depends upon the
ontological role it plays in the existent under
examination.

An invariant may be necessary for identity preservation
within a specified domain if its constitutive
role has been established. It may also contribute
to a sufficient set of conditions when considered
alongside other ontological relations.

Nevertheless, no invariant should be treated as
universally necessary or sufficient without an
independent argument.

The distinction between formal invariance and
ontological invariance must also be maintained.
A formally preserved relation cannot automatically
be regarded as an ontologically preserved condition
of identity.

\subsection{Manifestation and the Conditions of Identity Preservation}

Manifestations may provide relevant indications
of the conditions under which an existent persists
through transformation.

The continuity or alteration of manifestations
may reveal changes in the determinations and
relations through which an existent is expressed.

However, manifestation continuity cannot be treated
as a sufficient condition of identity preservation.

A sequence of similar manifestations may belong
to numerically distinct existents. Conversely,
the manifestations associated with one existent
may undergo substantial alteration without
independently establishing the termination of
its identity.

The continuity of manifestation must therefore
be distinguished from the preservation of identity.

The same distinction applies to the capacity
to manifest. The preservation of such a capacity
may be relevant to the investigation, but it
does not independently establish numerical identity.

The conditions of identity preservation must
therefore be examined without identifying identity
with the continuity of its manifestations or
the uninterrupted exercise of its capacities.

\subsection{Causal Relations and Identity Preservation}

Causal relations may contribute to an account
of identity preservation by establishing connections
between successive conditions of an existent.

A sequence of transformations may be causally
connected, and such connections may help explain
how one condition gives rise to another.

Nevertheless, causal continuity does not independently
establish identity preservation.

Causal relations may connect numerically distinct
existents. A later existent may arise through
a causal process involving an earlier existent
without thereby being numerically identical with it.

Likewise, the absence of an established causal
connection does not independently settle the
question of identity discontinuity.

Causal relations must therefore be distinguished
from numerical identity.

Their ontological relevance depends upon whether
the causal relation is constitutive of the identity
under examination or merely describes a relation
between successive conditions.

\subsection{The Problem of Necessary Conditions}

A condition is necessary for identity preservation
within a specified domain when its absence excludes
the possibility of preserving the relevant identity
under the conditions being considered.

The identification of such conditions requires
an independent ontological argument.

A feature that frequently accompanies identity
preservation is not necessarily required for it.
A condition that is important in one class of
existents may be irrelevant or differently relevant
in another.

Consequently, no material, structural, functional,
relational, or manifestational condition should
be declared universally necessary merely because
it is associated with familiar cases of persistence.

The investigation must establish:

\begin{enumerate}
    \item the existent whose identity is under examination;
    \item the domain and transformations involved;
    \item the proposed necessary condition;
    \item the ontological reason why its absence would
    exclude identity preservation.
\end{enumerate}

Without these specifications, claims of necessity
remain insufficiently determined.

\subsection{The Problem of Sufficient Conditions}

A condition or set of conditions is sufficient
for identity preservation within a specified domain
when its presence provides adequate grounds for
establishing that the numerical identity of the
existent has remained the same.

Sufficiency must not be confused with correlation,
similarity, or the preservation of a frequently
observed feature.

A condition may be strongly associated with identity
preservation without being sufficient to establish it.
Even a combination of preserved features may fail
to distinguish one numerically identical existent
from another existent exhibiting the same configuration.

The identification of sufficient conditions therefore
requires an account of why the relevant conditions
exclude the possibility of numerical replacement
or the emergence of a distinct identity.

Such an account must be specific to the ontological
domain under examination.

It is possible that identity preservation depends
upon a conjunction of conditions rather than a
single feature. Yet this possibility must itself
be established and must not be treated as a
universal assumption.

\subsection{Conjunctive and Relational Conditions}

The conditions of identity preservation may be
conjunctive or relational.

A particular condition may be insufficient when
considered independently but may contribute to
a sufficient account when combined with other
conditions.

For example, material constitution, structural
organization, and relational configuration may
have different ontological roles. Their combined
preservation may be relevant in a particular
domain even when none of them independently
establishes identity preservation.

However, the mere accumulation of preserved
features does not automatically produce a sufficient
criterion.

A conjunction of conditions must explain why
their combined presence is ontologically decisive
for the identity under examination.

The relation among conditions must therefore
be established rather than inferred from their
number or descriptive completeness.

The investigation must also remain open to the
possibility that some conditions are dependent
upon others, that some are relevant only under
specific transformations, and that some may
have different significance across distinct
domains of existence.

\subsection{Identity Preservation Without Complete Feature Preservation}

Identity preservation must not be assumed to
require the preservation of every feature of
an existent.

An existent may undergo transformations in
its determinations, material constitution,
structural organization, relations, or manifestations.
The alteration of such features does not
automatically establish that its identity has
ceased.

The possibility of identity preservation through
incomplete feature preservation is therefore
ontologically significant.

However, this possibility must not be interpreted
as an unrestricted claim that identity persists
regardless of the transformations involved.

The absence of complete feature preservation
does not prove identity continuity, just as
the preservation of many features does not
independently establish it.

The question remains whether the conditions
that have changed are compatible with the
persistence of the identity under examination.

This requires distinguishing between transformable
determinations and conditions whose alteration
would be incompatible with the continued identity
of the existent.

The distinction cannot be settled through
a universal list of preserved or transformable
features. It must be established in relation
to the ontological constitution of the existent
and the transformations involved.

\subsection{Identity Preservation Across Interruption}

Identity preservation may also be examined
under conditions of interruption.

An existent may undergo an interruption in
a particular process, manifestation, relation,
or mode of activity. Such interruption may
alter the conditions under which persistence
is described.

Nevertheless, the interruption of a particular
process or manifestation does not independently
establish the termination of identity.

Nor does the possibility of resumption establish
that the same identity has persisted throughout
the interruption.

The ontological relation between the conditions
before, during, and after interruption must
be examined independently.

The conditions of identity preservation must
therefore account for interruption without
equating it either with identity termination
or with an insignificant temporary absence.

The relevance of interruption depends upon
the particular existent, the interrupted
condition, and the respect under which identity
preservation is being examined.

\subsection{The Possibility of Multiple Identity Relations}

The conditions of identity preservation must
not presuppose that every existent is associated
with exactly one identity.

The relation between an existent and its identity
or identities remains an ontological question.
The preservation of one identity relation does
not automatically settle whether other identity
relations may also be involved.

Similarly, the multiplicity of components,
determinations, manifestations, or relations
does not independently establish a plurality
of identities.

The conditions of identity preservation must
therefore remain neutral concerning the number
of identities until the relevant ontological
relations have been established.

This neutrality is necessary to prevent a
particular model of unity or plurality from
determining the conclusions in advance.

\subsection{The Subject, Domain, and Respect of Identity Preservation}

Every claim concerning identity preservation
must specify its subject, domain, and respect.

The subject identifies the existent or identity
whose preservation is being examined.

The domain identifies the relevant class of
transformations, conditions, or temporal
succession.

The respect identifies the ontological aspect
under which preservation is asserted or denied.

A condition may be necessary in one domain
without being necessary in another. A feature
may be sufficient under a restricted class
of transformations without providing a universal
criterion.

Likewise, the preservation of identity in one
respect must not be generalized into a claim
of complete preservation in every respect.

These specifications are not merely methodological
conveniences. They are necessary for determining
the ontological scope and validity of claims
concerning identity preservation.

\subsection{Methodological Requirements}

The investigation of the conditions of identity
preservation must observe the following requirements:

\begin{enumerate}
    \item Identity preservation must be distinguished
    from the preservation of individual features.
    \item No material, structural, functional,
    relational, or manifestational condition
    should be assumed universally necessary
    or sufficient.
    \item The ontological relevance of every
    proposed condition must be independently
    established.
    \item Necessary conditions must be distinguished
    from sufficient conditions and from conditions
    that merely accompany persistence.
    \item Conjunctive conditions must be justified
    through their ontological relations rather
    than their descriptive accumulation.
    \item Invariance and manifestation continuity
    must not be treated as automatic criteria
    of identity preservation.
    \item Identity preservation through transformation
    must remain distinct from unrestricted persistence.
    \item The possibility of interruption must be
    examined without presupposing identity termination
    or uninterrupted preservation.
    \item The number of identities involved must
    remain open to independent investigation.
    \item Every claim must specify its subject,
    domain, and respect.
\end{enumerate}

These requirements establish a framework for
investigating identity preservation without
reducing it to any single form of continuity
or feature preservation.

\subsection{Scope of the Present Section}

This section has examined the ontological conditions
under which identity preservation may be established.

Material constitution, structural organization,
determination, relational configuration, invariance,
manifestation, and causal connection may contribute
to an account of identity preservation. Nevertheless,
none of these conditions can be treated as a
universal criterion without an independent
ontological justification.

The distinction between necessary and sufficient
conditions is essential. A feature that accompanies
identity preservation is not necessarily required
for it, and a feature that is required in a
particular domain may not be sufficient to establish
identity preservation.

The preservation of identity must therefore be
understood through an account of the ontological
relations that make the continued numerical identity
of a particular existent possible under specified
transformations.

The next section examines the limits of identity
continuity. It investigates the boundaries of
what can be established about identity preservation,
the conditions under which continuity claims remain
underdetermined, and the limits of extending
particular criteria across different domains
of existence.

% ============================================================
% End of Section 05
% ============================================================

% ============================================================
% Section 06 — Limits of Identity Continuity
% Book I: Ontology of Viability
% Chapter 05: Identity and Continuity
%
% Purpose:
%   To examine the ontological and methodological limits
%   of claims concerning identity continuity.
%
% Conceptual Policy:
%   1. Identity continuity must not be assumed through
%      unrestricted temporal persistence.
%   2. The limits of establishing continuity must be
%      distinguished from the limits of identity itself.
%   3. Insufficient grounds for continuity do not establish
%      non-identity.
%   4. Similarity, invariance, persistence, and causal
%      connection are not independently sufficient criteria.
%   5. No universal criterion of identity continuity may
%      be imposed without ontological justification.
%   6. Every claim must specify its subject, domain,
%      and respect.
% ============================================================

\section{Limits of Identity Continuity}
\label{sec:limits-of-identity-continuity}

\subsection{The Ontological Problem of the Limits of Continuity}

The investigation of identity continuity requires an
examination not only of the conditions under which
identity may be preserved, but also of the limits
within which such preservation can be established.

The possibility of transformation without the immediate
termination of identity does not imply that identity
must persist through every transformation. Likewise,
the failure to establish identity continuity under
particular conditions does not automatically establish
that identity has ceased.

The problem therefore concerns two distinct questions:
the ontological limits of identity continuity and the
limits of the grounds through which continuity can
be established.

The first concerns whether there are transformations
or conditions incompatible with the preservation
of a particular identity. The second concerns whether
the available ontological analysis provides adequate
grounds for determining whether identity has persisted.

These questions must not be conflated.

An inability to establish continuity may result from
an unresolved ontological relation, an inadequately
specified domain, or the absence of a sufficient
account of the relevant transformations. Such
underdetermination does not itself constitute proof
of discontinuity.

Conversely, the absence of an established contradiction
between transformation and identity preservation
does not establish unrestricted continuity.

The present section examines these limits while
maintaining the distinction between what is
ontologically possible and what has been adequately
established within a particular inquiry.

\subsection{The Limits of Temporal Generalization}

Identity continuity cannot be assumed to hold under
every temporal transformation.

The fact that an existent is considered across
successive temporal conditions does not establish
that its identity necessarily remains invariant
throughout the entire succession.

Temporal succession may involve transformations
that preserve identity, transformations that alter
conditions relevant to identity, or transitions
whose ontological significance remains unresolved.

Consequently, a general proposition asserting that
identity persists under all temporal transformations
requires an independent justification.

The temporal character of a transformation does
not determine whether that transformation is
compatible with identity preservation.

Nor can the continuity of time itself be used
as a sufficient ground for the continuity of
identity. Temporal continuity and identity continuity
are distinct ontological questions.

The scope of any claim concerning identity continuity
must therefore be restricted to the transformations
and conditions for which its validity has been
established.

\subsection{The Limits of Persistence as a Criterion}

Persistence cannot provide an unrestricted criterion
of identity continuity.

An existent may persist through a specified interval
without the persistence of its numerical identity
having been independently established. Similarly,
the interruption of a process or manifestation
does not automatically establish that identity
has ceased.

The persistence of an existent, the persistence
of identity, the continuity of manifestation,
and the continuity of the capacity to manifest
must therefore remain distinct.

A claim concerning one of these forms of persistence
cannot be transferred automatically to the others.

The limits of persistence as a criterion become
particularly significant when the subject of the
claim is left unspecified. The statement that
something persists may refer to a process, a
structure, a material configuration, a relation,
or an existent.

Without identifying the subject and respect,
persistence cannot provide a determinate conclusion
about identity continuity.

\subsection{The Limits of Invariance as a Criterion}

Invariance also has limited explanatory force
in relation to identity continuity.

A feature may remain invariant under a specified
transformation without establishing that the
numerical identity of its bearer has persisted.

The same invariant may characterize distinct
existents, and formal equivalence may obtain
between configurations that are not numerically
identical.

Conversely, the absence of an identified invariant
does not independently establish identity
discontinuity.

The limits of invariance arise from the distinction
between the preservation of a feature and the
preservation of identity.

An invariant can contribute to an account of
identity continuity only when its ontological
relevance has been established for the particular
existent and transformation under examination.

It cannot be elevated into a universal criterion
merely because it remains stable across a selected
class of transformations.

\subsection{The Limits of Similarity and Difference}

Similarity and difference provide no unrestricted
basis for determining identity continuity.

A high degree of similarity between successive
conditions does not independently establish
that they belong to one and the same numerical
identity.

Likewise, extensive qualitative difference does
not independently establish that the earlier
identity has ceased and been replaced by another.

The magnitude of similarity or difference may
be relevant to the investigation, but it does
not determine the ontological conclusion without
an account of the conditions under which the
comparison is made.

The same difficulty applies to gradual transformation.
A succession of small differences does not
necessarily establish continuity, just as the
accumulation of differences does not necessarily
establish discontinuity.

The limits of similarity and difference therefore
require that qualitative comparison remain distinct
from numerical identity.

\subsection{The Limits of Causal Continuity}

Causal continuity may establish a relation between
successive conditions, but it cannot independently
determine whether identity has persisted.

A causal sequence may connect distinct existents.
The emergence of a later existent from an earlier
one does not, by itself, establish that both
conditions belong to one and the same numerical
identity.

Likewise, an interruption in an identifiable
causal sequence does not automatically establish
the termination of identity.

Causal relations may be relevant to the conditions
of identity preservation in a particular domain.
However, their ontological significance must be
established rather than inferred from the existence
of causal connection alone.

The continuity of causal production must therefore
remain distinct from the continuity of identity.

\subsection{The Limits of Material and Structural Criteria}

Material and structural criteria may be relevant
to identity continuity, but their applicability
cannot be generalized without qualification.

Material preservation may be important for some
existents, while material replacement may be
compatible with the persistence of identity in
other cases.

Structural preservation may contribute to continuity
within a specified domain, but equivalent structures
may also occur in numerically distinct existents.

Likewise, structural transformation may occur
without independently establishing identity
discontinuity.

The ontological relevance of material and structural
conditions depends upon the particular existent,
its mode of constitution, and the transformations
under examination.

A criterion established for one class of existents
must not be extended to another class without
an independent argument.

\subsection{The Limits of Functional and Relational Criteria}

Functional and relational criteria also have
domain-dependent significance.

An existent may retain its function while its
identity is replaced by another existent capable
of performing the same function.

Conversely, an existent may undergo functional
transformation without necessarily losing its
identity.

The same applies to relational configurations.
A relational position may be occupied by distinct
existents, and one existent may participate in
different relational configurations across successive
conditions.

Functional and relational continuity therefore
cannot serve as universal criteria of identity
preservation.

Their significance must be established in relation
to the ontological constitution of the existent
and the specific domain under examination.

\subsection{The Limits of Manifestation-Based Continuity}

Manifestation-based accounts of continuity encounter
a further limitation.

The continuity of manifestations does not necessarily
establish the continuity of the identity associated
with them. Similar manifestations may belong to
distinct existents, and a continuous manifestation
sequence may not be sufficient to determine the
numerical identity of its source or bearer.

Conversely, an interruption or alteration in
manifestation does not independently establish
the termination of identity.

The same distinction applies to the continuity
of the capacity to manifest. The persistence of
such a capacity may be relevant, but it cannot
replace an ontological account of identity.

Consequently, manifestation-based criteria must
remain limited to the respect in which manifestation
continuity is actually established.

They cannot be extended automatically to claims
concerning the numerical identity of an existent.

\subsection{The Limits of Extrapolation Across Domains}

A criterion of identity continuity established
within one domain cannot automatically be generalized
to every other domain.

Different existents may possess different modes
of constitution, determination, organization,
and manifestation. A condition that is necessary
for identity preservation in one domain may be
non-essential in another.

Likewise, a sufficient condition established
under a restricted class of transformations
may cease to be sufficient when the transformation
class is expanded.

The validity of a continuity criterion therefore
depends upon the scope of the domain in which
it has been established.

This limitation prevents the construction of
universal ontological conclusions from particular
cases without adequate justification.

Any extension of a criterion beyond its original
domain requires an independent argument concerning
the relevant similarities, differences, and
ontological conditions.

\subsection{The Limits of Establishing Identity Discontinuity}

The establishment of identity discontinuity
requires more than the identification of change,
interruption, or the failure of a particular
continuity criterion.

A claim that an identity has ceased to persist
must establish that the relevant transformation
or condition is incompatible with the continued
numerical identity of the existent.

The absence of a previously identified invariant,
the loss of a material constituent, the alteration
of a structure, or the interruption of a process
does not independently provide such an argument.

The termination of identity must therefore be
distinguished from the failure to establish
its continuity.

This distinction does not imply that identity
can never cease or that every transformation
is compatible with its preservation. It requires
only that non-identity be established through
adequate ontological grounds rather than inferred
from the absence of a preferred form of continuity.

\subsection{The Limits of Establishing Identity Continuity}

The establishment of identity continuity also
has limits.

The preservation of selected features, the
existence of causal relations, the continuity
of manifestations, or the persistence of an
existent may provide relevant grounds, but they
do not necessarily establish numerical identity.

A continuity claim must explain why the successive
conditions belong to one and the same existent
rather than to distinct existents related by
similarity, reproduction, causation, or structural
equivalence.

Where this explanation is unavailable, the claim
of identity continuity remains underdetermined.

The absence of adequate grounds must not be
transformed into an affirmative assertion of
continuity merely because no decisive discontinuity
has been established.

Identity continuity and identity discontinuity
are both ontological claims requiring justification.

\subsection{Underdetermination and the Suspension of Judgment}

An important consequence of the preceding limits
is the possibility of ontological underdetermination.

A particular transformation may be established,
and its effects on material, structural, functional,
or relational conditions may be described, while
the question of numerical identity remains unresolved.

In such a case, neither continuity nor discontinuity
should be asserted without adequate grounds.

The suspension of judgment does not imply that
identity itself is indeterminate. It indicates
that the available ontological analysis has not
yet established which conclusion is justified.

The distinction between an unresolved claim and
an unresolved reality must therefore be maintained.

An ontological condition may be determinate even
when the inquiry has not adequately established
its nature.

Likewise, the absence of a justified conclusion
must not be interpreted as evidence that identity
is merely subjective or dependent upon recognition.

The limits of an inquiry do not automatically
constitute limits of the existent being investigated.

\subsection{The Problem of Irreversibility and Restoration}

Irreversible transformation and restoration introduce
additional questions concerning identity continuity.

An existent may undergo a transformation that
cannot be reversed in the relevant domain. The
irreversibility of that transformation does not
automatically establish identity discontinuity.

Conversely, the restoration of a previous configuration
does not independently establish the restoration
of the same numerical identity.

A later configuration may resemble an earlier
one without being numerically identical with
the existent that previously exhibited it.

Likewise, an irreversible change may occur within
the persistence of one identity, depending upon
the ontological conditions of the existent.

Irreversibility and restoration must therefore
be examined independently of the question of
numerical identity.

The relevant issue is not simply whether a previous
state can be recovered, but whether the identity
whose continuity is under examination has been
preserved through the transformation.

\subsection{The Limits of Unity and Plurality in Continuity}

The question of identity continuity must remain
open in relation to the possibility of multiple
identities.

The continuity of an existent does not independently
establish that exactly one identity is involved.
Likewise, the multiplicity of manifestations,
components, relations, or transformations does
not establish that multiple identities are involved.

A succession may concern the continuity of one
identity, the relation between distinct identities,
or a more complex ontological configuration.

These possibilities cannot be resolved by assuming
either the unity or the plurality of identity
in advance.

The number and relation of identities must therefore
be investigated independently of the mere continuity
or discontinuity of particular manifestations
and configurations.

\subsection{The Subject, Domain, and Respect of the Limits}

The limits of identity continuity must be specified
relative to the subject, domain, and respect of
the claim.

The subject identifies the existent or identity
whose continuity is being examined.

The domain identifies the transformations, temporal
conditions, or ontological relations within which
the claim is made.

The respect identifies the particular aspect
under which continuity or discontinuity is asserted.

A limitation established for the continuity
of manifestation cannot automatically be transferred
to the continuity of identity. A restriction
concerning material preservation cannot automatically
be generalized to every structural or relational
domain.

The scope of every conclusion must therefore
remain proportionate to the conditions under
which it has been established.

\subsection{Methodological Requirements}

The investigation of the limits of identity
continuity must observe the following requirements:

\begin{enumerate}
    \item Identity continuity must not be assumed
    under unrestricted temporal transformation.
    \item Persistence, invariance, similarity,
    causal continuity, and manifestation continuity
    must not be treated as universally sufficient
    criteria.
    \item The absence of adequate grounds for
    continuity must not be treated as proof
    of discontinuity.
    \item The absence of established discontinuity
    must not be treated as proof of continuity.
    \item Criteria established within a particular
    domain must not be generalized without
    independent justification.
    \item The termination of identity must be
    distinguished from the interruption of
    a process or manifestation.
    \item The restoration of a configuration
    must not be treated as sufficient proof
    of the restoration of numerical identity.
    \item The possibility of multiple identities
    must remain open to investigation.
    \item Underdetermination must be distinguished
    from ontological indeterminacy.
    \item Every claim must specify its subject,
    domain, and respect.
\end{enumerate}

These requirements establish the limits within
which claims concerning identity continuity
may be responsibly formulated.

\subsection{Scope of the Present Section}

This section has examined the limits of identity
continuity and the conditions under which claims
of preservation or discontinuity may remain
undetermined.

No single criterion based on persistence,
invariance, similarity, material constitution,
structural organization, function, relation,
causation, or manifestation can be treated as
universally sufficient for establishing identity
continuity.

Likewise, the failure to establish continuity
through a particular criterion does not independently
establish that identity has ceased.

The ontological limits of identity continuity
must therefore be distinguished from the limits
of a particular inquiry. Both continuity and
discontinuity require adequate grounds, and
neither may be inferred merely from the absence
of evidence for the other.

The final section synthesizes the distinctions
and arguments developed throughout this chapter.
It establishes the conceptual results of the
investigation while preserving the unresolved
questions concerning the ontological conditions
and possible forms of identity continuity.

% ============================================================
% End of Section 06
% ============================================================

% =============================================================================
% Section 07 — Synthesis of Identity and Continuity
% Book I — Ontology of Viability
% Chapter 05 — Identity and Continuity
%
% Purpose:
%   To synthesize the ontological distinctions and arguments developed
%   throughout the chapter, establish the conceptual status of identity
%   continuity, and clarify what can and cannot be concluded about
%   identity preservation across transformation.
%
% Conceptual Policy:
%   1. Identity is not reducible to persistence, invariance, or continuity.
%   2. Persistence of an existent and persistence of identity are distinct.
%   3. Continuity of manifestation is not equivalent to continuity of identity.
%   4. Qualitative difference does not establish numerical non-identity.
%   5. Identity preservation requires ontological argument, not merely
%      the observation of preserved features.
%   6. The unity or plurality of identities remains an open question.
%   7. The inquiry is ontological and independent of recognition.
% =============================================================================

\section{Synthesis of Identity and Continuity}
\label{sec:synthesis-of-identity-and-continuity}

\subsection{The Ontological Status of Identity Continuity}

The preceding sections have examined the relation between identity and continuity
through several distinct but interconnected problems: persistence, invariance,
transformation, difference, preservation, and the limits of continuity claims.
Their combined result is not a single criterion by which identity continuity can
be universally determined. Rather, it is the clarification of a set of
ontological distinctions that must remain intact whenever the continuity of
identity is investigated.

Identity continuity concerns the persistence of an identity through relevant
conditions of temporal succession, transformation, and possible interruption.
It does not designate a mere sequence of successive states, the preservation of
a collection of properties, or the uninterrupted continuation of manifestation.
Nor does it follow automatically from the persistence of the existent to which
an identity is attributed.

The question is therefore not simply whether something continues to exist, whether
its properties remain recognizable, or whether a continuous sequence of states
can be established. The question is whether the identity in question is
ontologically preserved across the transformations and conditions under
consideration.

This formulation requires a distinction between the existence of an existent
and the identity of that existent. Although existence and identity are
ontologically inseparable, their analytical examination is not identical.
Consequently, the persistence of an existent cannot be treated as an exhaustive
account of the persistence of its identity. Likewise, the termination of one
mode of manifestation does not, by itself, establish the termination of the
identity associated with it.

Identity continuity must therefore be investigated as a distinct ontological
problem rather than derived from a general principle of temporal persistence.

\subsection{The Distinction Between Identity, Persistence, and Continuity}

The analysis of the chapter establishes that several concepts commonly brought
together in discussions of persistence must be distinguished.

First, \emph{identity} concerns the ontological identity of an existent. It is
not equivalent to the continued presence of a particular property, structure,
material constitution, relation, or manifestation.

Second, \emph{persistence of an existent} concerns the continuation of that
existent through relevant conditions. This continuation may involve substantial
changes in its material, structural, relational, or manifestational
determinations.

Third, \emph{persistence of identity} concerns whether the identity itself is
preserved. It cannot be established solely by showing that the existent
continues, because the ontological relation between the existent and its identity
must itself be taken into account.

Fourth, \emph{temporal continuity} concerns the relation between conditions or
states across temporal succession. A continuous temporal sequence may provide
relevant conditions for investigating identity continuity, but temporal
continuity and identity continuity are not equivalent.

Fifth, \emph{continuity of manifestation} concerns the continuation of
manifestational conditions. Manifestation may continue, be interrupted, change,
or cease without these facts alone determining whether identity is preserved.

Finally, \emph{continuity of the capacity to manifest} concerns the persistence
of a capacity or possibility of manifestation. This must be distinguished both
from the actual continuation of manifestation and from the persistence of
identity.

These distinctions prevent the substitution of one kind of continuity for
another. They also prevent the assumption that every interruption in
manifestation constitutes an interruption in identity, or that every
uninterrupted sequence of manifestations establishes the persistence of one
and the same identity.

The relations among these concepts require investigation in each relevant
ontological domain and respect.

\subsection{Invariance as a Limited Condition of Identity Continuity}

The examination of invariance has demonstrated that invariance and identity
continuity are related but non-equivalent concepts.

An invariant is established relative to a transformation, a domain, and a
respect. Its preservation indicates that something remains unchanged under
the specified transformation in the specified respect. It does not, without
further argument, establish that the identity of an existent has been
preserved.

The preservation of an invariant may be relevant to identity continuity when
the invariant has an independently established ontological relation to the
identity under consideration. However, this relation cannot be assumed merely
because a feature remains constant.

Conversely, the failure to identify a preserved invariant does not establish
the discontinuity of identity. Identity may remain ontologically preserved
even where the examined properties, structures, material configurations, or
relations undergo extensive transformation.

Accordingly, neither of the following propositions can be accepted as a
universal principle:

\[
\text{Identity Continuity}
\equiv
\text{Invariance}
\]

\[
\text{Identity Discontinuity}
\equiv
\text{Absence of Identified Invariants}
\]

The stronger proposition that identity is invariant under all temporal
transformations is likewise unwarranted unless the relevant transformations
have first been shown to belong to a class under which identity is necessarily
preserved.

Temporal transformation, as such, does not establish the ontological
admissibility of every possible transformation with respect to identity.
The classification of transformations as identity-preserving, identity-altering,
or identity-terminating requires independent justification.

Invariance may therefore function as a relevant condition, an indicator, or
a component of a broader account of preservation. It cannot be elevated into
a universal definition or sufficient criterion of identity continuity.

\subsection{Difference, Transformation, and Numerical Identity}

The chapter has also established that qualitative difference and numerical
non-identity must not be conflated.

An existent may undergo material, structural, functional, relational, or
manifestational changes without these changes necessarily establishing that
it has become numerically different from what it was.

Similarly, the persistence of similarities between two conditions does not
establish that they belong to one and the same identity. Similarity may occur
between numerically distinct existents, while substantial qualitative
difference may occur within the transformations of an existent whose identity
is preserved.

The relevant distinctions can therefore be stated as follows:

\[
\text{Qualitative Difference}
\not\Rightarrow
\text{Numerical Non-Identity}
\]

\[
\text{Qualitative Similarity}
\not\Rightarrow
\text{Numerical Identity}
\]

Transformation must also be distinguished from replacement, termination,
and restoration. A transformation is not necessarily a replacement; a
replacement is not simply another name for transformation; and the
reappearance of a previous configuration does not, by itself, establish
the restoration of the same identity.

These distinctions are especially important where transformation occurs
gradually. No universal degree of qualitative difference can be assumed to
constitute a threshold at which numerical identity necessarily ceases.
Nor can the absence of a clearly identifiable discontinuity be treated as
proof that identity has remained continuous.

The ontological status of the relation between successive conditions must
therefore be established rather than inferred from the magnitude, direction,
or apparent reversibility of change alone.

\subsection{Conditions of Identity Preservation}

The investigation of preservation has shown that identity continuity may
depend on several kinds of ontological conditions. These may include material,
structural, determinative, relational, causal, and manifestational conditions,
as well as combinations of them.

No single category can be presumed to provide a universal criterion of
identity preservation.

Material continuity may be relevant in certain domains without being
necessary in every domain. Structural continuity may contribute to the
preservation of identity without being sufficient to establish it.
Determinative conditions may constrain the possibilities of transformation,
while relational conditions may be relevant to the persistence of an identity
within a particular ontological configuration.

Causal continuity may establish a significant relation between successive
conditions, but causal succession alone does not establish numerical identity.
Manifestational continuity may provide evidence of an ongoing ontological
process without exhausting the question of whether identity itself is
preserved.

The possibility of conjunctive conditions must also remain open. Identity
preservation may depend not on one isolated feature but on a configuration
of mutually related conditions. Such a configuration may involve different
conditions with different ontological roles and degrees of relevance.

The distinction between necessary and sufficient conditions must consequently
be maintained.

A condition shown to be necessary within a specified domain is not thereby
sufficient. Likewise, a condition shown to be sufficient under a restricted
set of assumptions cannot automatically be generalized to other domains.

Every claim about necessary or sufficient conditions must therefore specify
the relevant existent, the identity relation under consideration, the domain
of application, and the respect in which preservation is being asserted.

Identity preservation is not established by accumulating unchanged features.
It requires an account of why the preserved conditions are ontologically
relevant to the identity whose continuity is at issue.

\subsection{Manifestation, Interruption, and the Persistence of Identity}

The distinction between identity and manifestation has particular importance
for the problem of continuity.

Manifestation may continue without interruption, undergo transformation,
become intermittent, or cease under particular conditions. None of these
possibilities can be identified immediately with the preservation or
termination of identity.

The continuity of manifestation does not necessarily entail the continuity
of identity. A sequence of manifestations may exhibit continuity while
the question of whether they belong to one identity remains unresolved.

Conversely, an interruption in manifestation does not necessarily establish
the termination of identity. The absence of a manifestation, considered
in isolation, does not determine the ontological status of the identity
associated with it.

The continuity of the capacity to manifest introduces a further distinction.
An existent may retain a capacity for manifestation even when manifestation
is not actualized. However, the persistence of such a capacity does not
automatically establish the persistence of identity. Its ontological
significance must also be established.

The following distinctions must therefore remain explicit:

\[
\text{Persistence of Identity}
\neq
\text{Continuity of Manifestation}
\]

\[
\text{Continuity of Manifestation}
\neq
\text{Continuity of Capacity to Manifest}
\]

\[
\text{Continuity of Capacity to Manifest}
\neq
\text{Persistence of Identity}
\]

These distinctions allow the possibility of identity preservation across
manifestational interruption to be considered without presupposing that
every interruption is harmless to identity.

They also prevent the opposite assumption: that identity must cease whenever
its manifestation becomes discontinuous.

The relation between identity and manifestation must be determined through
the relevant ontological conditions rather than through the mere presence
or absence of manifestational activity.

\subsection{The Limits of Establishing Continuity and Discontinuity}

The preceding analysis does not establish that identity continuity is
always determinable through available ontological criteria. It establishes
that the question must be approached without reducing identity to any
single observable or formally specified condition.

There may be cases in which the conditions relevant to identity preservation
are sufficiently established to support a continuity claim. There may also
be cases in which the conditions establish termination or replacement.
Between these possibilities, however, there may remain cases in which the
ontological relation cannot be determined from the conditions presently
specified.

Such underdetermination must not be confused with either continuity or
discontinuity.

The absence of sufficient grounds for asserting continuity does not prove
discontinuity. Equally, the absence of sufficient grounds for asserting
discontinuity does not prove continuity.

This distinction is important because an ontological question may remain
underdetermined by a particular analysis without thereby becoming
ontologically indeterminate in itself. The limits of a given account,
classification, or formal model must not be mistaken for the limits of
being.

The inquiry must therefore distinguish between:

\[
\begin{aligned}
&\text{Ontological Underdetermination} \\
&\text{Analytical Insufficiency} \\
&\text{Unestablished Continuity} \\
&\text{Unestablished Discontinuity}
\end{aligned}
\]

These expressions designate different limitations and cannot be used
interchangeably.

The investigation of identity continuity is not grounded in the act of
recognizing an existent, remembering its previous conditions, or comparing
successive appearances from the standpoint of an observer. Such activities
may belong to other forms of inquiry, but they do not constitute the
ontological foundation of continuity in the present work.

The object of investigation is the ontological relation itself, whether
or not it is recognized, represented, or successfully determined by an
external inquiry.

\subsection{Identity Continuity and the Question of Unity and Plurality}

The question of identity continuity also remains connected to the problem
of unity and plurality examined in the preceding chapter.

The continuity of an existent does not, by itself, establish that the
existent possesses exactly one identity. Likewise, the plurality of its
manifestations, components, relations, or successive conditions does not
establish that it possesses multiple identities.

Unity of an existent, unity of identity, unity of source, and unity of
aggregate are distinct ontological questions. They cannot be treated as
equivalent simply because they may coincide in particular cases.

The possibility of multiple identities associated with one existent, as
well as the possibility of one identity across multiple manifestations
or configurations, must remain open to investigation.

Consequently, identity continuity cannot be formulated in advance as the
continuity of exactly one identity through every transformation of an
existent. Nor can changes in composition or manifestation be taken as
sufficient grounds for positing the emergence of a new identity.

The number and relation of identities must be established independently
of assumptions about the unity or plurality of manifestations and
components.

This preserves the distinction between the continuity of an existent,
the continuity of an identity, and the possible plurality of identities
associated with an ontological configuration.

\subsection{A Consolidated Ontological Position}

The results of this chapter can be consolidated into a set of propositions
that delimit the meaning and scope of identity continuity.

\begin{enumerate}

    \item Identity is not reducible to persistence, temporal continuity,
    invariance, similarity, or continuity of manifestation.

    \item The persistence of an existent and the persistence of its identity
    are distinct analytical questions, although existence and identity
    remain ontologically inseparable.

    \item Temporal succession does not by itself establish identity
    continuity, and temporal interruption does not by itself establish
    identity discontinuity.

    \item Invariance is relative to a transformation, domain, and respect.
    Its preservation establishes identity continuity only when its
    ontological relevance to the identity in question has been independently
    justified.

    \item The absence of an identified invariant does not establish the
    absence or termination of identity.

    \item Qualitative difference does not necessarily establish numerical
    non-identity, and qualitative similarity does not establish numerical
    identity.

    \item Transformation, replacement, termination, and restoration are
    distinct ontological possibilities and must not be conflated.

    \item Identity preservation may involve material, structural,
    determinative, relational, causal, manifestational, or conjunctive
    conditions. No single condition can be presumed universally necessary
    or sufficient.

    \item The continuity of manifestation and the continuity of the
    capacity to manifest are distinct from the persistence of identity.

    \item Interruption or absence of manifestation does not, by itself,
    establish identity termination.

    \item Continuity and discontinuity are both claims requiring
    ontological justification. The failure to establish either claim
    cannot be converted into proof of the opposite claim.

    \item The continuity of an existent does not settle the number of
    identities associated with it. Unity and plurality remain open
    ontological questions.

    \item Identity continuity is an ontological relation independent
    of recognition, memory, observation, or the successful representation
    of successive conditions.

\end{enumerate}

These propositions do not constitute a universal algorithm for deciding
every case of identity continuity. Their function is to establish the
conceptual and ontological constraints within which such decisions must
be justified.

They also prevent the question of identity preservation from being
resolved through circular definitions. Identity cannot first be defined
as an invariant and then declared persistent because an invariant has
been preserved. Nor can continuity be defined as the uninterrupted
persistence of identity and subsequently used as independent proof of
that persistence.

The ontological relation under investigation must remain distinct from
the criteria and conditions through which it is examined.

\subsection{The Scope of the Established Results}

The chapter has established a framework for examining identity continuity
without reducing it to any single temporal, structural, material, relational,
or manifestational criterion.

It has clarified the difference between persistence of an existent and
persistence of identity, between invariance and continuity, between
qualitative difference and numerical non-identity, and between continuity
of manifestation and continuity of identity.

It has also established that identity preservation must be investigated
through conditions whose ontological relevance is independently justified.
The necessary and sufficient status of such conditions remains domain-dependent
and cannot be generalized without argument.

Several questions nevertheless remain open.

The precise ontological grounds on which identity is preserved across
different classes of transformation have not been universally determined.
The possibility of continuity across particular forms of interruption
requires further specification. The relation between identity preservation
and the different forms of ontological determination remains subject to
further analysis. Likewise, the possibility and structure of multiple
identities cannot be resolved merely through the continuity of an existent.

These open questions do not invalidate the distinctions established here.
Rather, they indicate the limits of the conclusions that can legitimately
be drawn from the present analysis.

The central result is therefore methodological and ontological at once:
identity continuity is a real problem concerning the persistence of
identity, but it cannot be settled by substituting for identity any one
of the conditions with which it may be associated.

\subsection{Concluding Statement}

Identity continuity is neither the mere persistence of an existent nor
the preservation of an invariant feature through time. It is the
ontological question of whether an identity remains preserved across
the conditions, transformations, and possible interruptions to which
the relevant existent is subject.

The persistence of an existent, the persistence of identity, temporal
continuity, continuity of manifestation, continuity of the capacity
to manifest, invariance, and identity preservation are related but
non-identical determinations. Their relations must be established
without reducing one to another.

Likewise, qualitative difference cannot be treated as numerical
non-identity, and similarity cannot be treated as proof of numerical
identity. Neither continuity nor discontinuity may be inferred merely
from the absence of evidence for its opposite.

The inquiry therefore concludes without adopting a universal formula
for identity preservation. What has been established is the necessity
of examining identity continuity through ontologically relevant
conditions, while maintaining the distinction between identity itself
and the various forms of persistence, transformation, and manifestation
through which its continuity may be investigated.

Identity continuity is thus retained as a distinct ontological problem:
one whose resolution requires more than the preservation of features,
the continuation of processes, or the formal correspondence of successive
states.

% ============================================================
% END OF CHAPTER 05
% ============================================================